\documentclass{zapiski}
\originfo{}{}{}{2026}
\setcounter{page}{1}
\date{24 July 2026}

\usepackage[utf8]{inputenc}
\usepackage[T2A]{fontenc}
\usepackage[russian, english]{babel}

\usepackage{times}
\usepackage{url}
\usepackage{latexsym}
\usepackage{multirow}
\usepackage{graphicx}
\usepackage{amssymb}
\usepackage{amsmath}
\usepackage{amsthm}
\usepackage{hyperref}
\usepackage{booktabs}
\usepackage{tabularx}
\usepackage{caption}
\usepackage{cite}
\usepackage{paralist}
\usepackage{subcaption}
\usepackage{xcolor,colortbl}
\usepackage{array}
\usepackage{setspace}
\usepackage{placeins}
\usepackage{tikz}
\usetikzlibrary{arrows.meta}
\usepackage{pgfplots}
\pgfplotsset{compat=1.16}
\definecolor{ccal}{HTML}{009E73}
\pgfplotsset{
  qpanel/.style={width=.49\linewidth, height=3.4cm, grid=major,
    grid style={dashed,gray!14}, xmin=-35, xmax=95, ymin=-55, ymax=38,
    tick label style={font=\tiny},
    every axis label/.append style={font=\scriptsize},
    title style={font=\scriptsize\bfseries, yshift=-2pt}},
  bgpt/.style={only marks, mark=*, mark size=0.6pt, gray!30},
  stillocw/.style={only marks, mark=*, mark size=1.5pt, cadd},
  stillcr/.style={only marks, mark=*, mark size=1.5pt, cclamp},
  arocw/.style={-{Stealth[length=4pt]}, cadd!88!black, line width=1pt,
                quiver={u=\thisrow{u}, v=\thisrow{v}}},
  arcr/.style={-{Stealth[length=4pt]}, cclamp!88!black, line width=1pt,
               quiver={u=\thisrow{u}, v=\thisrow{v}}}
}
\newcommand{\qdecor}{%
  \fill[cgate!11] (axis cs:-35,-10.08) rectangle (axis cs:95,38);
  \draw[cgate!65!black, dashed, thick] (axis cs:-35,-10.08) -- (axis cs:95,-10.08);
  \node[anchor=north east, font=\tiny\itshape, text=cgate!60!black] at (axis cs:93,36) {gate ON};
  \node[anchor=south east, font=\tiny\itshape, text=gray] at (axis cs:93,-53) {gate OFF};
  \draw[ccal!75!black, thick, dash dot, rotate around={150:(axis cs:31,-4.6)}]
    (axis cs:31,-4.6) ellipse [x radius=30.5, y radius=13.7];
  \node[font=\tiny\itshape, text=ccal!60!black, anchor=west] at (axis cs:-33,-30) {calibrated zone};
  \addplot[bgpt] table[x=x,y=y,col sep=comma]{pic/viz/plane_mmlu.csv};
}

% pack floats tightly (the class default wastes half-pages around floats)
\renewcommand{\topfraction}{.95}
\renewcommand{\bottomfraction}{.9}
\renewcommand{\textfraction}{.05}
\renewcommand{\floatpagefraction}{.85}
\setcounter{topnumber}{3}
\setcounter{totalnumber}{5}
\captionsetup{font=small}
\setlength{\textfloatsep}{9pt plus 2pt minus 3pt}
\setlength{\intextsep}{9pt plus 2pt minus 3pt}
\setlength{\abovecaptionskip}{5pt}
\emergencystretch=1.8em

\pgfplotsset{
  planepanel/.style={width=.52\linewidth, height=4.6cm, grid=major,
    grid style={dashed,gray!15}, xmin=-35, xmax=95, ymin=-65, ymax=45,
    tick label style={font=\tiny}, every axis label/.append style={font=\scriptsize},
    title style={font=\scriptsize\bfseries, yshift=-2pt}},
  ptocw/.style={only marks, mark=*, mark size=1.1pt, cadd},
  ptcr/.style={only marks, mark=*, mark size=1.1pt, cclamp},
  ptrest/.style={only marks, mark=*, mark size=0.7pt, gray!55},
  disp/.style={-stealth, cgate, thick, quiver={u=\thisrow{u}, v=\thisrow{v}}}
}


\newtheorem{theorem}{Theorem}
\newtheorem{corollary}{Corollary}
\newtheorem{proposition}{Proposition}
\newtheorem{definition}{Definition}
\newtheorem{remark}{Remark}

\renewcommand{\UrlFont}{\ttfamily\small}

% figure palette (Okabe-Ito subset, CVD-validated; identity carried by
% direct labels as well, never by color alone)
\definecolor{cadd}{HTML}{E69F00}   % additive baseline
\definecolor{cclamp}{HTML}{0072B2} % clamp
\definecolor{cot}{HTML}{009E73}    % optimal transport
\definecolor{cgate}{HTML}{CC79A7}  % gated / NP-BW

\newcommand{\ocw}{OCW}
\newcommand{\crok}{CR}
\newcommand{\Wtwo}{W_2}

\english

\begin{document}

\title[Projection-Transport Steering]{Projection-Transport Steering: Conditional Control of Reasoning Behavior in LLMs}

% NOTE(team): all authors must be registered on OpenReview and added to the submission.
\author[Kashirskiy]{Mark Kashirskiy}
\address[M.~Kashirskiy]{ITMO University}
\email{353056@niuitmo.ru}

\author[Sergeev]{Daniil Sergeev}
\address[D.~Sergeev]{NUST MISIS}
\email{m2300826@edu.misis.ru}

\author[Ter-Hovhannisyan]{Tatevik Ter-Hovhannisyan}
\address[T.~Ter-Hovhannisyan]{}
\email{tatevik.terhovh@gmail.com}

% Mentor (per submission requirements the mentor must be listed):
\author[Piontkovskaya]{Irina Piontkovskaya (mentor)}
\address[I.~Piontkovskaya]{Huawei Noah's Ark Lab}
\email{irapiont@yandex.ru}

\begin{abstract}
Activation steering promises behavior control without retraining, but a global
shift of the residual stream suppresses useful and harmful behavior together.
In this work we treat steering as an optimal-transport problem on a
low-dimensional behavioral subspace and separate the decision to intervene
from the intervention itself. We prove that the transport map between
projection laws is the unique minimum-displacement intervention local to that
subspace, that the selectivity of any transfer function on an axis is capped
by how well the axis separates the subpopulations, and that a
Neyman--Pearson gate composed with the map is jointly optimal in its class. A
sequential form of the gate decides from a prefix, so the operator runs in a
single pass. On reasoning confidence in \texttt{gemma-2-2b-it} the additive
edit removes every overconfident-wrong answer only by destroying the
confident-right ones, which the token-level Cohen's $d$ of $0.049$ we measure
between the two makes a geometric necessity rather than a tuning failure.
Across five resampled MMLU subsets our conditional operator reaches
selectivity $+0.279$ at no accuracy cost, transfers to ARC-Challenge and
GSM8K without refitting, and replicates at $4.5\times$ scale. Giving every
baseline an equal tuning budget, we find it keeps that selectivity when the
questions change, where a tuned full-space map loses two thirds of its own. It
also leaves capability intact where the additive edit does not, staying at the
unsteered value on HumanEval, WikiText-103 and open-ended generation, and a
held-out OLMo-2/TruthfulQA evaluation supports the detector axis beyond Gemma.
\end{abstract}

\keywords{activation steering \and optimal transport \and calibration \and overconfidence \and mechanistic interpretability \and inference-time intervention.}

\maketitle

\section{Introduction}
\label{sec:intro}

Inference-time activation steering promises behavior control without retraining:
add a direction to the residual stream and the model refuses less, hedges more,
or changes style. The workhorse is contrastive activation addition
(CAA)~\cite{rimsky2024caa}: a fixed vector, a hand-tuned coefficient, applied to
every token. Yet applied to a concrete reasoning behavior --- overconfidence
of \texttt{gemma-2-2b-it} on multiple-choice reasoning --- CAA fails in a way
no coefficient tuning fixes: the dose that eliminates
\emph{overconfident-wrong} answers (\ocw{} $50\!\to\!0$ under a calibrated
readout) simultaneously eliminates \emph{confident-right} ones (\crok{}
$76\!\to\!1$) and costs $8.7$ accuracy points.\looseness=-1

The failure is not that additive steering is a weak tool. On refusal it is a
strong one, and prior work models that behavior as a single-direction
shift~\cite{arditi2024refusal,rimsky2024caa}, and a shift is exactly what an
additive edit implements optimally. Overconfidence asks for something else. The
target is not a relocated distribution but a compressed one: the tail must move
while the bulk stays. Corollary~\ref{cor:additive} covers both cases with one
statement, and reading it in the right direction is what separates the
successes from this collapse.

This paper treats that failure as a mathematical object rather than an
engineering nuisance. Our contributions are as follows. First, we give a
unifying framework with optimality theorems. Every linear steering operator is a scalar transfer
function $g$ on the projection $p=\langle h,\hat v\rangle$, with additive
steering $g(p){=}p{+}\alpha$, ablation $g{=}0$, and CAST-style
gating~\cite{lee2024cast} $g(p){=}p{+}\alpha\mathbf{1}[p{>}\tau]$. We prove in
Thm.~\ref{thm:ot} that the optimal-transport map between projection laws is the
unique minimal-perturbation intervention, and in Cor.~\ref{cor:additive} that
additive steering is optimal only for pure location shifts. Second, we give a
quantified impossibility result and its repair. Thm.~\ref{thm:gate}b caps the
selectivity of \emph{any} transfer function on an axis by the separability of
the subpopulations along it, and we measure a token-level Cohen's $d=0.049$
between \ocw{} and \crok{} on the confidence axis, so the additive collapse is
a geometric necessity rather than a tuning failure. The repair is conditional.
An LDA gate on a second, detection axis is Neyman--Pearson optimal, and
gate$\times$transport is jointly optimal in its class (Thm.~\ref{thm:joint}),
while Thm.~\ref{thm:seq} ties an online decision to the post-hoc one at a
stated rate and so keeps the whole pipeline single-pass. Third, we obtain
strong trade-offs at minimal action cost in the evaluated setting. Across five
MMLU seeds our conditional operators reach selectivity $+0.22$--$0.28$ at zero
accuracy cost, against $+0.00$ for additive steering at its published dose and
$+0.20$ for prompting (Sec.~\ref{sec:results}). Given every method an equal
tuning budget they keep that selectivity when the questions change, where a
faithful reimplementation of MiMiC keeps half of its own, and they retain more
\crok{} than any baseline that matches them (Sec.~\ref{sec:results}). They
replicate at $4.5\times$ scale, transfer to ARC-Challenge and GSM8K with no
refitting (Sec.~\ref{ssec:ninetb}), and every operator is $O(d)$ per token
with 19~KB of state, the optimal order by Prop.~\ref{prop:compute}. Our
results are a first account of \emph{why} a steering edit damages what it was
not aimed at, and of how much of that damage the geometry forces rather than
the tuning.\looseness=-1

\section{Related Work}
\label{sec:related}

\textbf{Linear steering.} ActAdd~\cite{turner2023actadd} and
CAA~\cite{rimsky2024caa} add fixed diff-in-means vectors,
ITI~\cite{li2023iti} shifts probe-selected heads, RepE~\cite{zou2023repe}
generalizes reading and control vectors, directional ablation removes one
direction everywhere~\cite{arditi2024refusal}, and SAE clamping fixes a feature
value~\cite{templeton2024scaling}. All are special cases of our
transfer-function family, and all are input-independent edits.
AxBench~\cite{wu2025axbench} finds them weaker than prompting on
instruction-following, which is the concrete gap this family has to close.
CAST~\cite{lee2024cast} adds conditional \emph{gating}, but its action
remains additive, and SADI~\cite{wang2025sadi} and KL-adaptive
composition~\cite{scalena2024dynamic} adapt dose per input without changing
its shape. In this paper we take the shape of the edit as the object of study,
and derive locality, dose control and serving cost from one
optimality statement rather than fixing them by hand.\looseness=-1

\textbf{Distribution matching and optimality.} MiMiC~\cite{singh2024mimic}
derives the Gaussian-$\Wtwo$ affine map, applied unconditionally in the full
space, AcT~\cite{rodriguez2025act} fits per-neuron 1-D transport, and
LinEAS~\cite{rodriguez2025lineas} learns affine maps under a distributional
loss. COAST~\cite{nguyen2026coast} proves \emph{pointwise} optimality of a
covariance-weighted projection, which in our language is transport to a point
mass under the $M$ of Prop.~\ref{prop:metric}, and A-LQR
\cite{nguyen2026alqr} proves LQR-optimality of a control input. We prove a different statement in Theorem~\ref{thm:ot}, namely the optimal map
between \emph{distributions} under a locality constraint, together with the
conditional optimality and selectivity ceiling of
Thms.~\ref{thm:gate}--\ref{thm:joint} that no prior steering work states. PCHI~\cite{liu2026probe} trains
probe-conditioned head rescalings for the same behavior, and its additive
baseline corrects $\approx0\%$, independently confirming
Cor.~\ref{cor:additive}. That baseline is ungated, so its design leaves the
gated-additive cell empty and confounds gating with the multiplicative head
parameterisation, and our grid fills exactly that cell. It also contains a
probe-conditioned member of the same family, a full-$d$ logistic gate composed
with each action (Appendix~\ref{app:fulltables}). \cite{zhang2026wired} reports the dose--response
our per-sample dose exploits.

\textbf{Deciding during generation.} We are not the first to move the decision
inside the loop. CAST~\cite{lee2024cast} compares each token's hidden state
with a condition vector at a tuned threshold, DSAS~\cite{ferrando2026dsas}
computes context-dependent per-token scales and states the same
when-versus-how split we use, and GAPS~\cite{gaps2026} adds dimension-level
gates at $O(d)$ per token. A prediction probe read after a few generated
tokens can anticipate whether steering will succeed~\cite{steerable2026}, and
YOPO~\cite{luo2026yopo} removes a second clean pass by reconstructing the
unsteered residual before reading its direction. What none of these provides
is a decision with an error guarantee: our gate is a Neyman--Pearson test whose
online decision is tied to the post-hoc one at a stated rate
(Thm.~\ref{thm:seq}). Two findings in this literature shape our design. Reading
a direction off already-steered activations degrades it~\cite{luo2026yopo}, so
our reader sits above the action and the decision is a single commitment, after
which the reader's output is no longer used. And a feature that
\emph{detects} a behavior in finished text is a poor \emph{predictor} of it
from a prefix~\cite{fpcg2026}, so we do not reuse the whole-trace head online:
we refit it on the prefix budget the gate will actually have, and report the
budget--AUROC curve (Appendix~\ref{app:layers}).

\section{Method: Projection-Transport Steering}
\label{sec:method}

We write $h\in\mathbb{R}^d$ for the residual-stream activation at a fixed
decoder layer and let $V\in\mathbb{R}^{k\times d}$ have orthonormal rows
spanning a \emph{behavioral subspace}. The two axes that carry the behavior
are a confidence axis and a correctness axis, obtained by diff-in-means on the
model's own labeled reasoning traces; $k$ itself is a design parameter that we
select rather than fix. We write $q=Vh\in\mathbb{R}^k$, let $\mu$ be the law
of $h$ under natural generation, $\mu_V$ the induced law of $q$, and $\nu$ a
target law on $\mathbb{R}^k$.

\begin{definition}[admissible steering]
$T:\mathbb{R}^d\to\mathbb{R}^d$ is admissible for $(V,\nu)$ if
(i) $(I-V^\top V)\,T(h)=(I-V^\top V)\,h$ $\mu$-a.s.\ (\emph{locality}: the
orthogonal complement is untouched), and (ii) $(V\!\circ T)_{\#}\mu=\nu$, i.e.\ the pushforward of $\mu$ equals the target law
(\emph{goal}: steered projections follow the target law). The cost is
$C(T)=\mathbb{E}_\mu\lVert T(h)-h\rVert^2$.
\end{definition}

We regularize by that cost because large interventions are what empirically
destroy generation, as we show in Sec.~\ref{sec:results}, and because under a
quadratic functional metric $C$ bounds the KL drift of the next-token
distribution (Prop.~\ref{prop:metric}). The first result says that once the
subspace and the target law are fixed, there is nothing left to design.

\begin{theorem}[steering is optimal transport of the projection]
\label{thm:ot}
$\min_T C(T) = \Wtwo^2(\mu_V,\nu)$ over admissible $T$, attained by
$T^*(h)=h+V^\top\!\big(S(Vh)-Vh\big)$ where $S$ is the ($\mu_V$-a.s.\ unique) optimal
Monge map from $\mu_V$ to $\nu$. Any optimal $T$ satisfies $V T(h)=S(Vh)$
almost surely, so the optimal intervention is a deterministic function of the
projection alone and conditioning on $h_\perp$ cannot reduce the cost.
\end{theorem}

\begin{proof}[Proof sketch]
Locality collapses the cost to the projection, and $(Vh,VT(h))$ couples
$(\mu_V,\nu)$, so $C(T)\ge\Wtwo^2$. Brenier's map attains it. Full proof in
Appendix~\ref{app:proofs}.
\end{proof}

For $k=1$ the map $S=F_\nu^{-1}\!\circ F_{\mu_V}$ is the monotone quantile
map, and we implement it as a 41-knot piecewise-linear lookup. For Gaussian
marginals it is affine~\cite{knott1984optimal,gelbrich1990formula},
$S(q)=m_t+A(q-m_s)$ with
$A=\Sigma_s^{-1/2}\big(\Sigma_s^{1/2}\Sigma_t\Sigma_s^{1/2}\big)^{1/2}\Sigma_s^{-1/2}$.

\textbf{Choosing the subspace.} Theorem~\ref{thm:ot} holds for any orthonormal
$V$, so $k$ is a design parameter and not a fixed property of the behavior. Two
interpretable axes make the geometry legible and are what the gate needs, but
the action is free to use more. The natural larger choice is the subspace an
unconditional full-space map actually moves: the mean shift $m_t-m_s$ together
with the leading eigenvectors of $A-I$, where $A$ is the Bures--Wasserstein
matrix of the source and target moments. Retaining the interpretable axes and
then the first $k$ of those directions costs $\Theta(kd)$ per token instead of
$\Theta(d^2)$. At the far end $V=I$ recovers the full-space Gaussian map that
MiMiC~\cite{singh2024mimic} applies unconditionally, so one family runs from a
single interpretable axis to a dense $\Theta(d^2)$ operator and
Theorem~\ref{thm:joint} gates every point of it. We select $k$ on the tuning split like every other
hyperparameter, and Sec.~\ref{sec:results} reports what that buys, which on
this behavior is less than the construction suggests
(Appendix~\ref{app:parity}).

Reading Theorem~\ref{thm:ot} backwards tells us when the workhorse edit is
already the right one, and when it cannot be.

\begin{corollary}[additive steering is forced to be global]
\label{cor:additive}
Additive steering $g(p)=p+\alpha$ is the optimal map iff $\nu=\mu_V\!*\!\delta_\alpha$
(pure shift). If the target compresses the spread, which is what
``hedge only abnormal confidence'' requires ($\sigma_t<\sigma_s$), the optimal
slope is $\sigma_t/\sigma_s<1$: the tail must move more than the bulk. A
slope-1 map matched to the tail over-transports the bulk by the same constant,
which is the observed \crok{} collapse. Moreover its selectivity between two subpopulations
with equal projection laws is identically zero.
\end{corollary}

If no function of the confidence coordinate can separate the two
subpopulations, we need a second coordinate and a decision rule on it. The next
result says what the best such rule is, and how much it can buy.

\begin{theorem}[optimal gate and selectivity ceiling]
\label{thm:gate}
Let a gate fire iff a trace statistic $s>\tau$, with
$\mathrm{TPR}(\tau)=P(s>\tau\mid \ocw)$, $\mathrm{FPR}(\tau)=P(s>\tau\mid \crok)$.
(a) For any objective increasing in TPR and decreasing in FPR, the
likelihood-ratio test is optimal at every operating point (Neyman--Pearson).
Under equal-covariance Gaussian class conditionals it reduces to the linear
discriminant $w=\Sigma^{-1}(\mu_{\ocw}-\mu_{\lnot\ocw})$.
(b) For any statistic with concavified ROC area $\mathrm{AUC}$,
$\max_\tau[\mathrm{TPR}-\mathrm{FPR}] \le 2\,\mathrm{AUC}-1$.
\end{theorem}

Gate and action can then be designed separately without losing optimality.

\begin{theorem}[joint optimality of gate $\times$ transport]
\label{thm:joint}
Among interventions that are local to $V$, identity off the gate event
$\Gamma=\{s>\tau\}$, and achieve target law $\nu_\Gamma$ on $\Gamma$, the
cost-minimal intervention applies Theorem~\ref{thm:ot}'s map for the conditional
law $(\mu_V\mid\Gamma)\to\nu_\Gamma$ on $\Gamma$; its cost is
$P(\Gamma)\,\Wtwo^2(\mu_V\!\mid\!\Gamma,\ \nu_\Gamma)$.
\end{theorem}

Theorem~\ref{thm:joint} treats the gate as a switch. Nothing forces that: the
same cost functional, asked for the best \emph{dose} rather than the best
event, answers with a graded rule whose threshold is measured instead of
tuned.

\begin{proposition}[the optimal dose is affine in the gate posterior]
\label{prop:dose}
Let $T_\lambda(h)=h+\lambda\,V^\top\!\big(S(Vh)-Vh\big)$ interpolate between the
identity and the map of Thm.~\ref{thm:ot}, and let the dose $\lambda$ be any
$[0,1]$-valued function of the gate statistic $s$. Write $\pi(s)=P(\ocw\mid s)$,
$d(s)=\mathbb{E}[\lVert S(Vh)-Vh\rVert^2\mid s]$, and let the behavioral gain of
a dose be linear in it with class slopes $\beta_1$ on \ocw{} and $-\beta_0$ on
\crok{}. Then the dose maximising gain minus displacement priced at $1/2c$ is
\[
\lambda^*(s)=\Big[\tfrac{c(\beta_0+\beta_1)}{d(s)}\big(\pi(s)-\tau\big)\Big]_0^1,
\qquad \tau=\frac{\beta_0}{\beta_0+\beta_1},
\]
nondecreasing in $\pi$, zero exactly on $\{\pi\le\tau\}$, and converging
pointwise to $\mathbf{1}[\pi(s)>\tau]$ as $c/d\to\infty$. The hard gate is thus
the free-displacement limit of the graded operator, and it is strictly
suboptimal at any finite price of displacement whenever $\pi$ has no atom at
$\tau$.
\end{proposition}

\begin{proof}[Proof sketch]
The dose is $s$-measurable, so the objective splits as an integral over $s$ of a
strictly concave quadratic in $\lambda$; maximise it pointwise on $[0,1]$. Full
proof in Appendix~\ref{app:proofs}.
\end{proof}

Two things follow that the switch does not give us. The operating point $\tau$
is a ratio of two measured rates, the collateral and the removal a full dose
produces, so it is read off the tuning split rather than searched; and because
the linear discriminant of Thm.~\ref{thm:gate}a is the log-likelihood ratio,
$\pi(s)$ is $\sigma(s)$ up to the prior term, which we fit on the extraction
traces. One slope $\rho=c(\beta_0+\beta_1)/d$ is left to search, in place of a
threshold and an action.

Theorem~\ref{thm:joint} says nothing about \emph{when} $\Gamma$ must be
decided. Scoring a finished trace and regenerating it is one implementation.
Deciding while the trace is written is another, and it is the one a server can
afford. Write $z_t$ for the mean projection over the first $t$ generated
positions and $z_\infty$ for the completed-trace summary the post-hoc gate
uses.

\begin{theorem}[single-pass gating]
\label{thm:seq}
(a) Among all gates measurable with respect to the first $t$ positions, the
likelihood-ratio test on $z_t$ is optimal at every operating point, and under
equal-covariance Gaussian class conditionals it is the linear discriminant
fitted on $z_t$ --- Theorem~\ref{thm:gate} applied to the statistic the gate
actually has. Its ceiling is $2\,\mathrm{AUC}(z_t)-1$, which is what a
prefix budget costs.
(b) If instead the gate insists on the whole-trace statistic, let the
contributions form a martingale difference sequence around $z_\infty$ with
conditionally $\sigma$-sub-Gaussian increments and fire at the first $t$ with
$s_t>\tau+B_t(\delta)$, where
$B_t(\delta)=1.7\,\sigma\sqrt{(\log\log 2t+0.72\log(5.2/\delta))/t}$ is the
stitched finite-law-of-the-iterated-logarithm boundary of
Howard et al.~\cite{howard2021confseq}.
Then it fires on a trace with $s_\infty\le\tau$ with probability at most
$\delta$, simultaneously over all $t$, and on a trace with margin
$\Delta=s_\infty-\tau>0$ it fires by the first $t$ with
$B_t(\delta)\le\Delta/2$, after $O(\sigma^2\Delta^{-2}\log\log(1/\delta))$
positions.
(c) Either way the composed operator is single-pass: the action starts at the
firing position, costs $O(d)$ per token, keeps $O(1)$ extra state, and
invalidates no cache entry, because positions written before the switch keep
the keys they were written with.
\end{theorem}

\begin{proof}[Proof sketch]
Part (a) is Theorem~\ref{thm:gate} with $z_t$ in place of $z$, part (b) is the
time-uniform confidence sequence of
Howard et al.~\cite{howard2020supermartingale,howard2021confseq}, and part (c)
follows from the order of the two hooks. We give the full proof in
Appendix~\ref{app:proofs}.
\end{proof}

\begin{remark}[detection and action need not share a depth]
\label{rem:depth}
Nothing in Theorems~\ref{thm:ot}--\ref{thm:seq} forces the gate to read the
layer the action writes. Placing the reader above the actor costs one position
of lag and lets each sit where it is best. On \texttt{gemma-2-2b-it} the
detection contrast peaks at layer~16, where the trace-level AUROC is $0.799$
against $0.769$ at layer~14, and the gap is larger for the prefix heads the
sequential gate uses ($0.740$ against $0.712$ at 16 generated tokens), so we
let the reader sit at layer~16 while the action stays at layer~14
(Appendix~\ref{app:layers}). The order also removes an interference the
single-pass setting would otherwise have: reading a direction off activations
the same operator has already edited degrades it~\cite{luo2026yopo}, and here
the decision is a single commitment taken before anything is edited.
\end{remark}

Two properties of the operator remain to be stated: that it does not depend
on the Euclidean metric, and that it is as cheap as any such operator can be.

\begin{proposition}[metric covariance]
\label{prop:metric}
If cost is $(\Delta h)^\top M(\Delta h)$ with $M\succ0$ (e.g.\ a Fisher/Gauss--Newton
pullback, so cost approximates the induced next-token KL), the admissible problem
reduces to optimal transport of $\mu_V$ to $\nu$ \emph{inside} $\mathbb{R}^k$
under the Mahalanobis cost $M_V=VMV^\top\succ0$: the optimal map is
$S_M=M_V^{-1/2}\!\circ\nabla\varphi\circ M_V^{1/2}$ with $\varphi$ convex, and the
operator is again $T^*(h)=h+V^\top(S_M(Vh)-Vh)$. Euclidean PTS is $M_V=I_k$.
Proof in Appendix~\ref{app:proofs}.
\end{proposition}

\begin{remark}[support preservation]
The quantile map never leaves the target's empirical support. An additive
map violates support for every $\alpha\neq0$, consistent with the observed
generation collapse at $|\alpha|=1.5\lVert h\rVert$.
\end{remark}

\begin{proposition}[computational optimality]
\label{prop:compute}
Any algorithm implementing a non-degenerate intervention local to a subspace
spanned by a dense direction must read all $d$ coordinates of $h$, i.e.\ costs
$\Omega(d)$ per token. PTS costs $\Theta(kd+\log K)$, which matches that lower
bound for constant $k$ and $K$, against $\Theta(d^2)$ for a full-space affine
map (MiMiC), $\Theta(d\,d_{\mathrm{sae}})$ for SAE steering, and a full extra
forward pass for KL-adaptive dosing. Proof in Appendix~\ref{app:proofs}.
\end{proposition}

\section{Experimental Setup}
\label{sec:exp}\label{sec:data}

\textbf{Models, datasets, and splits.} The main experiments use \texttt{gemma-2-2b-it}~\cite{gemma2024}, and the
central comparison is repeated on \texttt{gemma-2-9b-it} at $4.5\times$
scale (Sec.~\ref{ssec:ninetb}).\footnote{2B weights were served from the
\texttt{unsloth/gemma-2-2b-it} mirror. We verified all 288 tensors are
bit-identical to the official \texttt{google/gemma-2-2b-it} release.}
Both models reason inside \texttt{<think>...</think>} and answer
\verb|\boxed{letter}|. Decoding is greedy, so evaluation ``seeds''
resample \emph{questions}, not decoding noise. Activations are read and
modified at layer 14/26 (2B) and 24/42 (9B). Benchmarks:
MMLU~\cite{hendrycks2021mmlu} (fitting and primary evaluation),
ARC-Challenge~\cite{clark2018arc}, and GSM8K~\cite{cobbe2021gsm8k} as
4-option MCQ with seeded distractors (GSM8K-MCQ, Appendix~\ref{app:gsm}),
$n{=}300$ questions each, disjoint from all fitting data. We fit directions, maps, and gate parameters \emph{once} per scale, on a
400-question MMLU extraction split ($\sim$50k trace tokens), and reuse them
everywhere, so cross-benchmark results are transfer rather than refits. We run
the main comparison on five resampled MMLU subsets and the 9B replication on
three.

\textbf{Behavioral states and confidence readouts.} Every answer gets one of four labels --- \emph{overconfident-wrong} (\ocw),
\emph{nonconfident-wrong}, \emph{nonconfident-right},
\emph{confident-right} (\crok) --- under two independent readouts: the
\emph{letter-logit} readout M2 (softmax of the letter logits at the
trace's own answer position) and the \emph{yes-probability} readout M4
(normalized $P(\mathrm{YES})$ over per-option yes/no reformulations:
\texttt{"Answer whether the given answer is the right answer\ldots{}
Question: \{q\} Answer: \{option\}"}, one forward per option, reasoning in
\texttt{<think>}, then \texttt{YES}/\texttt{NO}, and the four $P(\mathrm{YES})$
values are renormalised). M3 is the model's own stated probability, elicited
after the trace with \texttt{"How confident are you in this answer? Reply with
a probability between 0 and 1."} On
identical questions they disagree sharply (ECE $0.39$ against $0.10$,
Appendix~\ref{sec:measure}): M2 is strongly overconfident, M4 much better
calibrated. M4 is therefore the primary readout, and agreement under M2 and
the M3 self-report is a robustness requirement. Full grids:
Appendix~\ref{app:fulltables}.\looseness=-1

\textbf{Directions and the trace-level gate.} On the extraction split we fit a \emph{confidence} axis $\hat v$, the
steering axis, and a \emph{detection} axis $\hat u$, both by diff-in-means, and
we orthonormalize them into $V=(\hat v,\hat u)$. We keep the two contrasts
disjoint by construction. The first separates the two \emph{confident} states
$\{\ocw,\crok\}$ from the two nonconfident ones, so $\hat v$ encodes confidence
irrespective of correctness, and the second separates \ocw{} from \crok{}, so
$\hat u$ encodes whether that confidence is warranted. Layer~14 is the depth at which
the second contrast is most linearly readable (Appendix~\ref{app:layers}). A completed trace of length $T$
is summarized by its mean projections
$z=\big(\tfrac1T\sum_t\langle h_t,\hat v\rangle,\
\tfrac1T\sum_t\langle h_t,\hat u\rangle\big)$, and an LDA head gives the gate
score $s(z)=w^\top z+b$, firing iff $s(z)>\tau$. Here $w$ and $b$ come from the
equal-covariance LDA solution and $\tau$ is a quantile of the extraction-split
score distribution (the median unless stated), never a quantity fitted on
evaluation data. We run the gate in two regimes. \emph{Post-hoc} (two-pass):
generate unsteered, keep that output if the gate is silent, and otherwise
regenerate with the intervention active during reasoning. \emph{Sequential} (single-pass,
Thm.~\ref{thm:seq}): a second LDA head, fitted on the mean projection over the
first $t^\star$ generated positions rather than over the finished trace, is
read at layer~16 while the action runs at layer~14 (Remark~\ref{rem:depth}),
and the action switches on at token $t^\star$ and stays on. Fitting the head on
the statistic the gate will actually have is what makes an early decision
sound, and $t^\star$ and the threshold quantile are searched on the tuning
split like every other hyperparameter. Nothing is regenerated. In both regimes the
gate sees only the model's own activations, with no gold labels at
inference.\looseness=-1

\textbf{Steering methods.} All operators act through $p_t=\langle h_t,\hat v\rangle$
(Sec.~\ref{sec:method}). \emph{Additive} is $h_t'=h_t+\alpha\hat v$ at
$\alpha=-0.75\lVert h\rVert$, the dose that eliminates \ocw{} under the
calibrated readout. \emph{Ablation} is $h_t'=h_t-p_t\hat v$, and \emph{clamp}
is $p_t'=\min(p_t,c)$ at a source quantile, with the median as the
representative condition. \emph{Quantile, Gaussian and Bures--Wasserstein
transport} map onto a target law fitted on the extraction split. We form \emph{conditional variants} by composing the trace gate with each
action, giving CAST-style gate$+$additive, gated clamp, gated ablation, gated
OT and LDA-gated subspace BW. In each of them the gate controls
\emph{exposure} and the action controls \emph{severity}, especially on false
positives. We compare against prompting, additive
CAA~\cite{rimsky2024caa}, ablation~\cite{arditi2024refusal},
CAST-style~\cite{lee2024cast}, MiMiC~\cite{singh2024mimic} and Linear-AcT
per-neuron transport~\cite{rodriguez2025act}, all reimplemented and run
on identical records, decoding, readouts, and metrics. Every method, ours
included, is given the same tuning budget of eight configurations, searched on
the same held-out MMLU tuning split under one pre-declared rule. CAST gets its
threshold and dose, MiMiC and Linear-AcT get layer selection, and our gate gets
the depth it reads and the quantile it fires at
(Appendix~\ref{app:parity}).
We also run \emph{think-only variants}, which disable steering before the
answer tokens.

\textbf{Metrics and statistical protocol.} We choose each metric to answer one question about the intervention.

\emph{Accuracy} asks whether it hurt task performance, and we report the
change $\Delta$acc against the unsteered model.

\emph{ECE}, the expected calibration error, asks whether the confidence
numbers can be trusted. We compute the mean absolute gap between average stated
confidence and actual accuracy over 10 confidence bins, where lower is better
and $0$ means confidence matches reality.

The \emph{state histogram} asks where the answers moved, and we report the
counts of the four behavioral states before and after
intervention.\looseness=-1

The \emph{surgical metrics} ask whether it hit only its target. Pairing
answers by question, $\mathrm{ocw\_rm}=P(\text{left \ocw}\mid\text{was \ocw})$ is
the fraction of the harmful state removed,
$\mathrm{cr\_keep}=P(\text{stayed \crok}\mid\text{was \crok})$ the
fraction of the healthy state preserved, and selectivity
$\mathrm{Sel}=\mathrm{ocw\_rm}-(1-\mathrm{cr\_keep})$ combines them
($+1$ is a perfect scalpel, $0$ is indiscriminate).\looseness=-1

\emph{Selective prediction}, which we measure by AURC and E-AURC, asks whether
confidence still ranks right answers above wrong ones. AURC is the area under
the risk--coverage curve, where lower is better. Monotone post-hoc recalibration leaves it
invariant, so only interventions that change answers or their ranking can
move it.

Single-subset results carry paired-bootstrap 95\% CIs, and multi-subset
results are mean$\pm$std over independently resampled subsets.\looseness=-1

\textbf{Post-review validation protocol.} We froze a simplified detector-axis
action after a development-only Qwen3 component analysis, and we then evaluated
it on held-out TruthfulQA panels for Mistral-7B, OLMo-2-7B, and Granite-8B. Each
eligible architecture used the same 237-question composition (50 in-domain,
187 shifted) and was compared with a norm-matched spherical action using
10,000 paired BCa resamples, and the three architecture comparisons used
Bonferroni-adjusted two-sided intervals. OLMo-2 subsequently faced seven
pre-frozen construction and specificity controls. A separate six-cell
Mistral/OLMo-2/Granite $\times$ BBQ/ETHICS observer screen kept fitting,
screening, and confirmatory units disjoint, and only Mistral/ETHICS advanced to a
1,000-scenario confirmatory test, with intervals adjusted for the six planned
cells. Only observer-eligible cells were opened, and the planned denominator remained six.

\section{Results on MMLU}
\label{sec:results}

\begin{figure}[!t]
\centering
\begin{tikzpicture}
\begin{axis}[width=.62\linewidth, height=2.7cm, xlabel={projection onto confidence axis},
  ylabel={CDF}, grid=major, grid style={dashed,gray!20},
  legend style={font=\scriptsize,draw=none, at={(1.03,0.5)}, anchor=west},
  every axis label/.append style={font=\footnotesize},
  tick label style={font=\scriptsize}, ymin=0, ymax=1]
\addplot[cadd, very thick] table[x=overconfident_wrong, y=q, col sep=comma]{pic/cdf_m4conf.csv};
\addlegendentry{overconfident-wrong}
\addplot[cclamp, very thick, dashed] table[x=confident_right, y=q, col sep=comma]{pic/cdf_m4conf.csv};
\addlegendentry{confident-right}
\addplot[cot, thick] table[x=nonconfident_right, y=q, col sep=comma]{pic/cdf_m4conf.csv};
\addlegendentry{nonconfident-right}
\addplot[cgate, thick, dashed] table[x=nonconfident_wrong, y=q, col sep=comma]{pic/cdf_m4conf.csv};
\addlegendentry{nonconfident-wrong}
\end{axis}
\end{tikzpicture}
\caption{Token-level projection CDFs on the confidence axis by behavioral
state (extraction split, $\sim$50k tokens). The \ocw{} and \crok{} curves are
nearly identical (Cohen's $d=0.049$), while the confident and nonconfident
pairs are well separated ($d\approx0.6$).}
\label{fig:cdf}
\end{figure}

\textbf{The confidence axis does not separate justified from unjustified
confidence.} We first test the geometric premise of selective steering. In
Figure~\ref{fig:cdf} the \ocw{} and \crok{} projection CDFs nearly coincide at
$d=0.049$, while the confident and nonconfident pairs separate at
$d\approx0.6$. We read this as the axis encoding \emph{whether} the model is
confident and not whether that confidence is warranted, so no transfer function
of this coordinate can be surgical (Cor.~\ref{cor:additive} and
Thm.~\ref{thm:gate}b at $\mathrm{AUC}\approx.51$). The measurements agree: the selected dose
removes all 50 \ocw{} but takes \crok{} from 76 to 1. Across five subsets
additive selectivity is $+0.000\pm0.000$, accuracy drops $8.7\pm5.3$
points, and AURC worsens from $0.404$ to $0.556$. We therefore do not read the removal of \ocw{} as evidence of calibration,
because the shift suppresses desirable and undesirable confidence
alike. The reliability literature reports the same quantity from the other
side: steerability varies by behavior rather than by prompt construction, and
what predicts it is how well the two subpopulations separate along the
difference-in-means line~\cite{tan2024generalisation,braun2025unreliability}.
Thm.~\ref{thm:gate}b is the statement that this is not a correlation to be
improved by better fitting but a ceiling.\looseness=-1

\textbf{Projection-dependent actions reduce collateral damage.}
Median clamping --- one scalar touching only the abnormal tail --- reaches
selectivity $+.124$ on the representative subset at $+1.3$ accuracy
points. Across five subsets accuracy is $+0.001\pm0.034$, that is
\emph{preserved} rather than improved (single-split accuracy gains should be
treated as noise until seed-replicated), while selectivity holds at
$+0.145\pm0.039$. Quantile transport to the non-\ocw{} target reaches
$+0.119\pm0.086$, likewise without the additive loss. We read the lesson as one about shape rather than strength. Abnormal
high-confidence activations must move more than ordinary ones, and ablation,
clamp, and transport achieve that by conditioning on the input
projection.\looseness=-1

\textbf{Main result.} Across five MMLU subsets we find that gated ablation attains
selectivity $\mathbf{+0.279\pm0.090}$ (additive $+0.000\pm0.000$,
prompting $+0.198\pm0.086$) at no detectable accuracy cost ($-0.003\pm0.050$),
with the best ECE ($.094\pm.030$) and AURC ($.395\pm.038$), beating
faithful reimplementations of CAST-style, MiMiC, and Linear-AcT at their
published configurations on every column
(Table~\ref{tab:headtohead}, Fig.~\ref{fig:pareto}).\looseness=-1

\begin{figure}[!t]
\centering
\begin{tikzpicture}
\begin{axis}[width=.9\linewidth, height=4.4cm, xlabel={selectivity},
  ylabel={$\Delta$ accuracy}, clip=false, grid=major, grid style={dashed,gray!20},
  every axis label/.append style={font=\footnotesize},
  tick label style={font=\scriptsize}, scaled ticks=false,
  yticklabel style={/pgf/number format/fixed},
  xticklabel style={/pgf/number format/fixed},
  xmin=-0.02, xmax=0.42, ymin=-0.10, ymax=0.10]
\addplot[only marks, mark=*, mark size=2.6pt, cadd] coordinates {(0.2321,0.0000)}
  node[right=3pt,font=\scriptsize]{additive \cite{rimsky2024caa}};
\addplot[only marks, mark=*, mark size=2.6pt, cadd] coordinates {(0.1232,0.0067)}
  node[right=3pt,font=\scriptsize]{CAST-style \cite{lee2024cast}};
\addplot[only marks, mark=*, mark size=2.6pt, cadd] coordinates {(0.0974,-0.0167)}
  node[below=3pt,font=\scriptsize]{prompting};
\addplot[only marks, mark=triangle*, mark size=3.0pt, cot] coordinates {(0.0700,0.0133)}
  node[above left=1pt,font=\scriptsize]{MiMiC \cite{singh2024mimic} ($\Theta(d^2)$)};
\addplot[only marks, mark=triangle*, mark size=3.0pt, cot] coordinates {(0.1816,0.0033)}
  node[above=3pt,font=\scriptsize]{Linear-AcT \cite{rodriguez2025act}};
\addplot[only marks, mark=square*, mark size=2.6pt, cclamp] coordinates {(0.1237,0.0133)}
  node[above right=1pt,font=\scriptsize]{clamp $q_{50}$};
\addplot[only marks, mark=square*, mark size=2.6pt, cclamp] coordinates {(-0.0642,-0.0267)}
  node[below right=1pt,font=\scriptsize]{gated $k$-dim transport};
\addplot[only marks, mark=diamond*, mark size=3.6pt, cgate] coordinates {(0.1811,0.0267)}
  node[above right=2pt,font=\scriptsize]{\textbf{gated full-rank transport (ours)}};
\addplot[only marks, mark=diamond*, mark size=3.6pt, cgate] coordinates {(0.2179,0.0433)}
  node[above left=2pt,font=\scriptsize]{\textbf{score-proportional dose (ours)}};
\addplot[only marks, mark=diamond*, mark size=3.6pt, cgate] coordinates {(0.2763,0.0267)}
  node[below=11pt,font=\scriptsize]{\textbf{gated ablation (ours)}};
\addplot[only marks, mark=pentagon*, mark size=3.6pt, cgate] coordinates {(0.0616,-0.0067)}
  node[above right=2pt,font=\scriptsize]{\textbf{sequential gate (ours)}};

\addplot[gray, dashed] coordinates {(-0.02,0) (0.42,0)};
\end{axis}
\end{tikzpicture}
\caption{The selectivity--accuracy frontier under the identical MMLU protocol
(M4, representative subset, with multi-subset means in
Tab.~\ref{tab:headtohead}). Additive-action methods (orange) pay accuracy for
selectivity or achieve neither, ungated distribution matching (green) is safe
but weakly selective, and gate $\times$ tail-local action (ours) sits above
both at $\Theta(kd)$ per token against $\Theta(d^2)$ for a full-space map
(Prop.~\ref{prop:compute}; at this scale that difference is below the
wall-clock noise floor, Appendix~\ref{app:serving}).}
\label{fig:pareto}
\end{figure}

\begin{table}[!t]\footnotesize\setlength{\tabcolsep}{4pt}
\centering
\resizebox{\linewidth}{!}{%
\begin{tabular}{l|ccccc}
\hline
method (M4, MMLU, seeds mean$\pm$std) & $\Delta$acc & ECE & AURC & cr\_keep & Sel \\
\hline
% seeds: 7,11,23,31,47 (n=300 each, M4 readout)
baseline & acc $0.473\pm0.026$ & $0.099\pm0.019$ & $0.404\pm0.029$ & --- & --- \\
prompting & $-0.005\pm0.020$ & $0.118\pm0.018$ & $0.392\pm0.023$ & $0.735\pm0.063$ & $+0.198\pm0.086$ \\
additive CAA \cite{rimsky2024caa} & $-0.004\pm0.041$ & $0.080\pm0.017$ & $0.400\pm0.038$ & $0.621\pm0.067$ & $+0.256\pm0.133$ \\
CAST-style \cite{lee2024cast} & $-0.007\pm0.032$ & $0.072\pm0.012$ & $0.396\pm0.017$ & $0.795\pm0.031$ & $+0.227\pm0.082$ \\
MiMiC \cite{singh2024mimic} & $-0.015\pm0.045$ & $0.113\pm0.010$ & $0.408\pm0.040$ & $0.763\pm0.031$ & $+0.132\pm0.103$ \\
Linear-AcT \cite{rodriguez2025act} & $+0.001\pm0.032$ & $0.101\pm0.011$ & $0.410\pm0.027$ & $0.868\pm0.048$ & $+0.146\pm0.141$ \\
gated local action (ours) & $-0.015\pm0.036$ & $0.089\pm0.013$ & $0.413\pm0.037$ & $0.713\pm0.050$ & $+0.206\pm0.059$ \\
gated full-rank transport (ours) & $+0.003\pm0.021$ & $0.108\pm0.022$ & $0.394\pm0.044$ & $0.854\pm0.038$ & $+0.146\pm0.042$ \\
score-proportional dose (ours) & $-0.005\pm0.054$ & $0.089\pm0.010$ & $0.406\pm0.050$ & $0.650\pm0.056$ & $+0.231\pm0.120$ \\
sequential gate (ours, single pass) & $+0.009\pm0.015$ & $0.111\pm0.018$ & $0.398\pm0.024$ & $0.897\pm0.017$ & $+0.079\pm0.050$ \\
\hline

\end{tabular}}
\caption{Multi-seed head-to-head at equal tuning budget: every method runs the
configuration it won on the held-out tuning split (Appendix~\ref{app:parity}),
over five evaluation subsets of $n{=}300$. Greedy decoding is deterministic, so
seeds resample \emph{questions}, a stricter robustness axis than sampling
seeds. Only Linear-AcT reports seed variance in this literature. The
default-configuration grid is Table~\ref{tab:defaults}.}
\label{tab:headtohead}
\end{table}

We find tuning changes what the baselines are. The additive dose we select
is a third of the published one and the CAST dose is half, and at those doses
both stop being destructive: additive reaches $+0.256\pm0.133$ and CAST
$+0.227\pm0.082$ at no accuracy cost. We read that as the useful result of
running the protocol rather than a threat to the framework, because what the
two buy their selectivity with is visible in the same row. Additive retains
$0.621$ of \crok{} and CAST removes only $0.432$ of the \ocw{} it fires on,
the two ways Cor.~\ref{cor:additive} says a shift can trade. MiMiC holds the
residual stream together but keeps half its tuning-split selectivity once the
questions change ($+0.132\pm0.103$) and rewrites the stream at $\Theta(d^2)$
per token, $1152\times$ the cost of a $k{=}2$ edit. Prompting reaches
$+0.198\pm0.086$ and Linear-AcT $+0.146\pm0.141$. Our gated local action
holds $+0.206\pm0.059$ with retention $0.713$ and the smallest spread in the
table, and the score-proportional dose reaches $+0.231\pm0.120$. We do not
claim per-metric maxima in every run; we claim the strongest \emph{joint}
trade-off across removal, retention, accuracy and selective prediction
(Fig.~\ref{fig:pareto}), and the only operators here that cost nothing at all
where the behavior is absent (Table~\ref{tab:transfer}).\looseness=-1

\textbf{The comparison survives equal tuning.} The first question to ask of a
head-to-head is whether the baselines were tuned as hard as the method, so we
gave every operator the same budget of eight configurations on the same
held-out tuning split under the same selection rule
(Appendix~\ref{app:parity}). Tuning moves the baselines, and in one case a
long way. The additive dose falls from $-0.75\lVert h\rVert$ to
$-0.25\lVert h\rVert$ and its tuning-split selectivity rises from
$\approx0$ to $+0.271$, which is worth stating plainly: the additive edit is
not useless, it was over-dosed, and the dose that removes every \ocw{} is far
past the dose that trades best. MiMiC reaches $+0.263$ at layer~18, CAST
$+0.184$ at the \crok{} 70th percentile and Linear-AcT $+0.184$ at layer~14.
Our gated ablation fires at the 30th rather than the 50th percentile for
$+0.245$, and the score-proportional dose of the next paragraph reaches
$+0.267$, within noise of the tuned additive dose and above every other method. Table~\ref{tab:headtohead} carries those
configurations, unchanged, to the five evaluation subsets.\looseness=-1

\textbf{The subspace dimension is not the lever.} Theorem~\ref{thm:joint} says
our gate composes with any transport, so a full-space map is not a rival
family but the $k{=}d$ member of ours, and we tuned the family end to end
under one budget. We measure selectivity peaking at $k{=}2$, $+0.186$, and flat above
it, $+0.148$ at $k{=}8$ and again at $k{=}32$, so the two interpretable axes
already carry this behavior and the extra $2302$ coordinates buy nothing. The
$k{=}d$ member says the same thing from the other end. We measure the gate composed with MiMiC's own map
at $+0.286$ on a $400$-question split and $+0.155$ on the $1200$-question one, the largest movement of any configuration we searched
(Table~\ref{tab:split}), which is what a configuration selected on noise looks
like when the noise is removed. We keep $k{=}2$ because the data says the
behavior lives there, not because the operator is cheaper there. The other
members of the family, a token-local action and a posterior-weighted dose,
are tuned and carried the same way in Appendix~\ref{app:ablation}.

\textbf{What survives a change of questions.} Carried unchanged to five
resampled question sets, the tuning-split ordering does not hold. MiMiC keeps
half of what the layer search bought it, $+0.263$ on the split and
$+0.132\pm0.103$ on the subsets, and gating the same map stabilises it at a
lower level, $+0.155$ and $+0.146\pm0.042$. Our gated local action moves from
$+0.245$ to $+0.206\pm0.059$ and the score-proportional dose from $+0.267$ to
$+0.231\pm0.120$. A $\Theta(d^2)$ map fitted at one layer is reading the
sample, and a two-dimensional edit behind a gate has little to overfit.

\textbf{Spending the edit where the evidence is.} What the tuning data points
at is retention rather than removal: at its selected threshold the hard gate
fires on most flagged traces and full ablation on the marginal ones is what
costs \crok{}. So we let the strength of the action follow the gate score
instead of switching on it, the interpolation $g_\lambda$ the method already
defines, and that is the best configuration any method reaches under the
budget, $+0.267$ against $+0.271$ for the tuned additive dose and $+0.263$ for
MiMiC. It does not win the way the diagnosis suggested. The
selected configuration widens exposure instead of narrowing it, firing at the
10th percentile and scaling each edit by its margin, so grading the dose is
what makes a permissive threshold affordable. On the evaluation subsets it
reaches $+0.231\pm0.120$, and on ARC-Challenge, off the split it was tuned on,
it is the best member we have at $+0.342$ with retention $0.815$, against
$+0.290$ for the hard switch.\looseness=-1

\textbf{Gate and action shape are complementary.}
We can see \emph{why} gate and local action compose in the measured movement of
real regenerated traces (Fig.~\ref{fig:plane}, Appendix~\ref{app:plane}). Additive steering moves all 50 \ocw{} far past the calibrated
zone and drags all 76 \crok{} along. Gated ablation protects twice: 30
\crok{} below
threshold stay \emph{bitwise identical} (77/300 questions overall), and
of the 46 \crok{} it wrongly flags, 25 keep their state because the
tail-local action barely moves typical activations. CAST-style has the
first protection but not the second, and ungated clamp or transport has the
second but not the first.\looseness=-1

\textbf{Detection is the main remaining bottleneck.} We find the LDA gate separates \ocw{} from \crok{} at trace-level AUROC
$\approx0.77$: it catches $41/50$, misses 9, and false-fires on $46/76$
\crok{}. The diagnostic value $2\,\mathrm{AUROC}-1\approx0.538$ is close to
twice the achieved selectivity. Thm.~\ref{thm:gate}b applies to the
\emph{concavified} ROC area, so raw AUROC is not an exact behavioral ceiling.
Detection, rather than the transport map, is nevertheless the binding
empirical constraint, and Fig.~\ref{fig:benchgeo} (bottom) plots the gap
against the label oracle. Five different actions inside the same gate buy
removal at the same rate in retention (Appendix~\ref{app:ablation}), which is
the practical form of that constraint. Because the constraint is detection, we ask where detection is best. We refit
the gate at nine depths (Appendix~\ref{app:layers}) and find the separation is
not monotone in depth. It peaks at layer~16, where AUROC is $0.799$ and the
ceiling $0.599$, and collapses by layer~25 to $0.649$, where the residual
stream has already committed to an answer. Since the framework never required
the gate to read the layer the action writes, we let it read layer~16 and keep
acting at layer~14.\looseness=-1

\textbf{The edit costs nothing outside the behavior it was fitted for.}
Selectivity inside MMLU says nothing about what an always-on edit does to
generation it was never meant to touch, so we measure three tasks the
directions never saw (Appendix~\ref{app:capability}). We find that the dose
that removes
every \ocw{} takes HumanEval pass@1 from $.396$ to $.128$, doubles WikiText-103
perplexity ($23.4\!\to\!47.3$), and makes $57.5\%$ of 200 open-ended
generations degenerate (distinct-3 $.953\!\to\!.480$). A judge asked to compare the steered and unsteered answers to the same
instruction, in both presentation orders, prefers the unsteered one on
$185$ of $200$ prompts. We find the projection-local actions invisible on all
four measurements even when left on for every token, at pass@1 $.390$--$.396$,
perplexity $23.0$--$23.4$, distinct-3 $.946$--$.951$, order-consistent win
rates of $.470$--$.508$ against a tie value of $.5$, and no degenerate
generation. MiMiC is close behind at pass@1 $.396$ and win rate $.512$, for
$\Theta(d^2)$ per token. The equal-budget dose is the fair comparison and it
tells the rest of the story: at $-0.25\lVert h\rVert$ the additive edit is no
longer destructive, at pass@1 $.408$, perplexity $23.6$, win rate $.495$ and
no degenerate generation, so what the parity protocol selects is a dose too
small to break the model. We still find it the wrong shape. That dose leaves $69.9\%$ of MMLU
tokens outside the $99\%$ conditional ellipsoid against $0.9\%$ unsteered and
$1.0$--$1.3\%$ for every projection-local operator
(Appendix~\ref{app:manifold}), and it cannot abstain: our sequential gate
fires on none of HumanEval and none of the teacher-forced WikiText, so in
those two rows we report the unsteered numbers to the digit. Because the target law leaves
the bulk in place, we expect a projection-local operator to have nothing to
do where the behavior is absent, and that is what we observe.\looseness=-1

\textbf{Robustness across instruments, and think-only steering.} Absolute
counts vary across readouts, but the conclusions persist: gated ablation
reaches selectivity $+0.35$ (M4), $+0.15$ (M2), $+0.17$ (M3 self-report:
$37\%$ of \ocw{} removed at $0.80$ retention). Think-only variants keep
every effect's direction at near-zero accuracy cost: the intervention
acts through the reasoning trajectory, not the answer
logits.\looseness=-1

\section{Scale, Transfer, and Serving Cost}
\label{ssec:ninetb}

\textbf{Replication at 9B.} We find the 2B story replicates on
\texttt{gemma-2-9b-it} (Table~\ref{tab:ninetb}): additive steering has
identically zero selectivity and \crok{}-retention on every subset while
losing $18.4\pm3.3$ points. CAST-style buys $+0.283$ selectivity for
$11.4\pm3.2$ points and $0.45$ retention, while the local operators alone gain
accuracy and \emph{improve} selective prediction (AURC $0.184$--$0.186$
against $0.197$), at selectivity $+0.20$ and retention $0.89$--$0.91$. We conclude that scale makes the global shift costlier rather than more
selective.\looseness=-1

\begin{table}[!t]\footnotesize\setlength{\tabcolsep}{4pt}
\centering
\resizebox{\linewidth}{!}{%
\begin{tabular}{l|ccccc}
\hline
method (M4, 9B, 3 seeds) & $\Delta$acc & ECE & cr\_keep & AURC & Sel \\
\hline
% 9B seeds [7, 11, 23]
baseline & --- & $0.144\pm0.021$ & $0.197\pm0.025$ & --- \\
additive $-0.75$ & $-0.184\pm0.033$ & $0.000\pm0.000$ & $0.450\pm0.023$ & $+0.000\pm0.000$ \\
CAST-style gate$+$add & $-0.114\pm0.032$ & $0.453\pm0.038$ & $0.285\pm0.016$ & $+0.283\pm0.093$ \\
clamp $q_{50}$ (ours) & $+0.023\pm0.012$ & $0.887\pm0.015$ & $0.184\pm0.025$ & $+0.202\pm0.020$ \\
gated ablation (ours) & $+0.011\pm0.008$ & $0.907\pm0.014$ & $0.186\pm0.023$ & $+0.203\pm0.034$ \\
\hline

\end{tabular}}
\caption{\texttt{gemma-2-9b-it}, MMLU, layer 24, $n{=}300\times3$ seeds:
the head-to-head replicates at $4.5\times$ scale.}
\label{tab:ninetb}
\end{table}


\textbf{Cross-domain transfer.} We apply the MMLU-fitted directions, maps, and
thresholds to ARC-Challenge and GSM8K-MCQ with no refitting
(Table~\ref{tab:transfer}), and we plot the two geometries in
Fig.~\ref{fig:benchgeo} (top).
\emph{ARC-Challenge} carries a real overconfident tail and the gated operators
transfer: our score-proportional dose reaches $+0.342$ at $\Delta$acc $-0.003$
with retention $0.815$, our gated local action $+0.290$ at $\Delta$acc
$+0.003$ with retention $0.901$, against $+0.184$ for CAST and $+0.151$ for
MiMiC. The tuned additive dose is higher still at $+0.385$, and keeps $0.801$.
\emph{GSM8K-MCQ} has only $6/300$ baseline \ocw{}: the tail targeted on MMLU
is absent, our gates read that correctly and fire on nothing, so every
conditional operator of ours is the exact identity there ($\Delta$acc
$0.000$, retention $1.000$), while the same tuned additive dose removes all
eleven \crok{} and the tuned CAST gate keeps three of them and costs $6.3$
accuracy points. An operator that can abstain pays nothing where the behavior
is absent, and an unconditional one cannot tell the difference. \emph{Threshold transfer} is the part of this that needs care. An absolute
MMLU threshold carried to another domain can fire on none of it or on all of
it, and which way it errs depends on the axis the threshold sits on. The
detector-axis gate at its default quantile over-fires on GSM8K-MCQ and takes
\crok{} with it, at retention $0.182$ in Table~\ref{tab:defaults}; the tuned
LDA gate errs the other way and fires on nothing, which is the direction we
want an unrefitted threshold to fail in. Label-free re-anchoring of $\tau$ to
a domain gate-score quantile removes the choice
(Appendix~\ref{app:fulltables}). The sequential gate needs no repair on this
domain either, because a gate that reads a prefix and commits once degenerates
to the identity where the behavior is absent, which is the behavior the
framework asks of it.\looseness=-1

\begin{table}[!t]\footnotesize\setlength{\tabcolsep}{3pt}
\centering
\resizebox{\linewidth}{!}{%
\begin{tabular}{l|ccc|ccc|ccc}
\hline
 & \multicolumn{3}{c|}{MMLU} & \multicolumn{3}{c|}{ARC-Challenge} & \multicolumn{3}{c}{GSM8K-MCQ} \\
method (M4, tuned) & $\Delta$acc & ECE & Sel & $\Delta$acc & ECE & Sel & $\Delta$acc & ECE & Sel \\
\hline
baseline & --- & 0.123 & --- & --- & 0.092 & --- & --- & 0.280 & --- \\
additive $-0.75$ & $-0.083^{*}$ & 0.099 & $+0.013$ & $-0.123^{*}$ & 0.265 & $+0.000$ & $-0.257^{*}$ & 0.069 & $+0.000$ \\
clamp $q_{50}$ & $+0.053^{*}$ & 0.116 & $+0.234^{*}$ & $+0.037$ & 0.116 & $+0.197^{*}$ & $+0.007$ & 0.286 & $-0.083$ \\
quantile-OT non-\ocw{} & $+0.013$ & 0.075 & $+0.158^{*}$ & $-0.027$ & 0.073 & $+0.095$ & $+0.010$ & 0.278 & $+0.083$ \\
gated clamp $q_{50}$ & $+0.037$ & 0.105 & $+0.248^{*}$ & $+0.027$ & 0.124 & $+0.137$ & $+0.020$ & 0.295 & $+0.167$ \\
\hline

\end{tabular}}
\caption{Cross-domain transfer with no refitting: every method keeps the
configuration it won on the MMLU tuning split and is applied to
ARC-Challenge and GSM8K-MCQ. $^{*}$ marks a bootstrap interval excluding
zero. GSM8K-MCQ has $6$ baseline \ocw{} and $11$ \crok{}, so its selectivity
column is noisy by construction, and the accuracy column is not.}
\label{tab:transfer}
\end{table}

\textbf{Transfer does not push activations off the manifold.} A map fitted on
one domain and applied to another is where we would expect an activation edit
to leave the region the model actually visits. We find it does not, provided
the edit is projection-local. Conditioning the edited coordinates on the
untouched complement gives a statistic that is $\chi^2_k$ under the null, so
the fraction of tokens outside the $99\%$ conditional ellipsoid has a value to
compare against rather than a threshold to tune. Unsteered traces sit at
$0.9$--$2.9\%$. The published additive shift puts $100\%$ of tokens outside it
on all three domains and the equal-budget dose still puts $40$--$70\%$ there,
while clamping, quantile transport, MiMiC and Linear-AcT stay at
$1.0$--$4.5\%$, that is at the null (Appendix~\ref{app:manifold}). The same
ordering appears in the nearest-neighbour distance to a bank of natural
layer-14 activations, where the additive shift lands $1.88\times$ further from
the model's own states than the states it replaces on MMLU and further still
off-domain, at $2.25\times$ on ARC and $2.04\times$ on GSM8K, and every
projection-local operator stays between $1.00\times$ and $1.28\times$. We
stress that the edit is small because the target law says it should be, not
because we tuned it to be.\looseness=-1

\textbf{Geometry and behavior disagree, and both matter.} The two diagnostics
do not rank the operators the same way, which is why we report both. Measured
by the KL the edit induces on the next-token distribution, the accepted
locality criterion in this literature~\cite{arditi2024refusal}, the worst
operator is a \emph{random} direction at the same norm as the additive dose:
$0.44$ on MMLU against a $0.1$ threshold, where the additive dose itself sits
at $0.032$ and every projection-local operator between $0.000$ and $0.006$. Yet
that random direction is invisible to the conditional test, because it barely
projects onto the behavioral coordinates the test watches. The conditional test
asks whether the behavioral coordinates were moved to an implausible place and
the KL asks whether the next token noticed. An operator has to pass both, and
only the projection-local ones do.\looseness=-1

\textbf{Post-review detector-axis validation.} We report the positive
predeclared held-out check in Table~\ref{tab:postreview}. On OLMo-2/TruthfulQA
the detector-axis action improves both primary metrics over a norm-matched
spherical action and passes all seven pre-frozen construction controls. We can
therefore rule out intervention magnitude alone as the explanation, and we find
that fitting the representation at prompted answer positions is necessary
relative to the unprompted ablation. We read this as evidence that the
detector-axis construction is not specific to Gemma.

\begin{table}[!t]\footnotesize\setlength{\tabcolsep}{3pt}
\centering
\resizebox{\linewidth}{!}{%
\begin{tabular}{llrrl}
\hline
model / panel & comparison & $\Delta$ calibrated MC2 [adj. 95\% CI] & $\Delta$ MC1 [adj. 95\% CI] & decision \\
\hline
OLMo-2 / TruthfulQA & matched spherical & $+.137\ [+.105,+.168]$ & $+.105\ [+.025,+.177]$ & pass \\
OLMo-2 / TruthfulQA & no intervention & $+.299\ [+.251,+.345]$ & $+.287\ [+.177,+.384]$ & pass \\
OLMo-2 / TruthfulQA & strongest isotropic & $+.220\ [+.180,+.258]$ & $+.194\ [+.097,+.287]$ & pass \\
OLMo-2 / TruthfulQA & unprompted fit & $+.295\ [+.247,+.343]$ & $+.287\ [+.177,+.384]$ & pass \\
\hline
\end{tabular}}
\caption{Held-out OLMo-2 validation on 237 TruthfulQA questions. Intervals use
10,000 paired BCa resamples and the predeclared Bonferroni family. The full
seven-control family passes, and the table shows the matched action and
representative construction controls. Qwen3 was development-only and is excluded from
confirmatory claims.}
\label{tab:postreview}
\end{table}

\textbf{Qualitative effects.} We find gated ablation does not induce generic
hedging: hedge markers per trace drop from $0.25$ to $0.13$ while traces get
shorter, from $86$ to $71$ words. The contrast with the additive dose is not
subtle. Under it only $28\%$ of traces still close their \texttt{<think>}
block and $20\%$ still emit an answer, distinct-3 falls from $0.91$ to $0.31$,
and $77\%$ of traces are degenerate. Under ours those numbers are unchanged
from the unsteered model (Table~\ref{tab:traces}). Samples are in
Appendix~\ref{app:samples}, and both lexicons are fixed in advance
(Appendix~\ref{app:defs}).\looseness=-1

\textbf{One pass, not two.} Theorem~\ref{thm:seq} lets the same gate decide
from a prefix, and the cost of that is measurable rather than assumed. Fitting
the LDA head on the first 16 generated positions instead of on the finished
trace costs $0.059$ of detector AUROC, $0.740$ against $0.799$
(Appendix~\ref{app:layers}), and the tuning split selects exactly that budget
with the threshold at the 30th percentile. The operator then takes one
decision per trace, at token 16, and regenerates nothing. Measured end to end
on the whole 300-question workload on one idle H100, that is the difference
between $45.5$ and $85.8$ seconds against $43.5$ unsteered, or $1.05\times$
against $1.97\times$, because the post-hoc gate must score every finished trace
and regenerate the $223$ it flags.

\textbf{Serving cost.}\label{sec:serving} Per token and per behavior, we spend one
$d$-dim dot product, a 41-knot lookup, and one FMA ($\approx4d$ FLOPs
$\approx10^{-4}$ of a layer, the optimal order of Prop.~\ref{prop:compute});
the gate adds a dot product and an $O(1)$ running mean, for a total state of
$\approx$19~KB. The action needs no extra forward pass and no KV-cache
invalidation. Honesty about what the wall clock does and does not show is
worth a sentence. At this scale it does not separate the single-pass
operators from each other: every one of them, MiMiC's $\Theta(d^2)$ map
included, lands within $6\%$ of the unsteered time, because the extra
arithmetic hides behind a memory-bound decode
(Appendix~\ref{app:serving}). What it separates is the number of passes. Strength interpolates as
$g_\lambda=(1-\lambda)\,\mathrm{id}+\lambda g$, and the operator matches the
vLLM residual-stream hook interface~\cite{vllmlens}.\looseness=-1

\section{Limitations and Assumptions}

The optimality results bound activation displacement under a stated subspace,
target law and locality constraint, with the usual regularity conditions on the
Monge map. They do not show the projection is sufficient for downstream
behavior, and the gate bound uses the concavified ROC, not raw AUROC.
Evaluation seeds resample questions under deterministic decoding, not training
runs. Architecture transfer is confirmed on OLMo-2/TruthfulQA, while other
architecture--behavior cells did not clear every predeclared gate
(Appendix~\ref{app:postreview}), so we claim no universal detector axis.

\section{Conclusion}

We turned an empirical steering failure into theorems and a method. The
additive edit is the wrong shape for a tail-compression target, optimal
transport along a behavioral subspace is the provably minimal intervention,
selectivity is constrained by detection as much as by the action, and the
NP-gate $\times$ tail-local action gives the strongest joint trade-off in our
Gemma/MMLU comparison across five subsets and two model scales, at no
detectable accuracy cost and minimal per-token compute. Two findings we did
not expect are worth carrying forward. Equal tuning rehabilitates the additive
baseline, and the subspace a full-space operator has over ours turns out not
to be where its advantage lives: on a tuning split large enough to identify a
configuration, selectivity peaks at $k{=}2$ and the gated full-rank member
falls to $+0.155$, so the comparison resolves inside our own framework rather
than against it. And the separation between
deciding and acting is what makes the method cheap to serve: the decision can
be taken on a prefix, with a stated agreement rate against the decision a
completed trace would have produced, so nothing is regenerated. We also find
that held-out validation on OLMo-2/TruthfulQA supports the detector-axis
construction beyond Gemma. We see three directions for further work: better
detection axes, transport under a Fisher metric, and further
behaviors.\looseness=-1

\section{Contributions}
% Per submission requirements: contribution of each team member.
\textbf{Mark Kashirskiy}: problem formulation of surgical steering; the PTS
framework, Theorems~1--4, Prop.~2 and proofs; the sequential gate and the
$k$-dimensional subspace construction; all steering operators,
baseline reimplementations, and experiments (2B and 9B, five seeds, four
benchmarks); visualizations; serving analysis; paper draft.
\textbf{Daniil Sergeev}: the confidence measurement suite (M1--M5) and its
$n{=}1000\times3$ study; feature-extraction pipeline (CAA/probe/SAE); the
original steering grid that exposed the global-hedging failure; the 9B
measurement replication; compute infrastructure.
\textbf{Tatevik Ter-Hovhannisyan}: the proposal of conditional (gated)
steering; behavioral labeling methodology and evaluation protocols (jointly
with D.S.: the LLM-as-judge and labeled-dataset design for the measurement
suite and the next behavioral domains); initial text review of the paper;
project direction and coordination.
\textbf{Irina Piontkovskaya} (mentor): project direction, research
supervision, methodology review; we thank her for guidance throughout
the project.

\providecommand{\bysame}{\leavevmode\hbox to3em{\hrulefill}\thinspace}
\begin{thebibliography}{10}

\bibitem{arditi2024refusal}
Andy Arditi, Oscar Obeso, Aaquib Syed, Daniel Paleka, Nina Panickssery, Wes
  Gurnee, and Neel Nanda, \emph{Refusal in language models is mediated by a
  single direction}, Advances in Neural Information Processing Systems 37, 2024.

\bibitem{clark2018arc}
Peter Clark, Isaac Cowhey, Oren Etzioni, Tushar Khot, Ashish Sabharwal, Carissa
  Schoenick, and Oyvind Tafjord, \emph{Think you have solved question answering?
  Try ARC, the AI2 reasoning challenge}, arXiv:1803.05457, 2018.

\bibitem{cobbe2021gsm8k}
Karl Cobbe, Vineet Kosaraju, Mohammad Bavarian, Mark Chen, Heewoo Jun, Lukasz
  Kaiser, et al., \emph{Training verifiers to solve math word problems},
  arXiv:2110.14168, 2021.

\bibitem{ferrando2026dsas}
Javier Ferrando, Xavier Suau, Pau Rodriguez, and Adri\`a Gonzalez,
  \emph{Dynamically scaled activation steering}, arXiv:2512.03661, 2026.

\bibitem{fpcg2026}
\emph{Detection features are not prediction features: forecasting model
  behavior from internal states}, arXiv:2606.11172, 2026.

\bibitem{gaps2026}
\emph{GAPS: gated activation steering at dimension granularity},
  arXiv:2609.01878, 2026.

\bibitem{gelbrich1990formula}
Matthias Gelbrich, \emph{On a formula for the $L^2$ Wasserstein metric between
  measures on Euclidean and Hilbert spaces}, Mathematische Nachrichten
  \textbf{147} (1990), no.~1, 185--203.

\bibitem{gemma2024}
Gemma Team, \emph{Gemma 2: Improving open language models at a practical size},
  arXiv:2408.00118, 2024.

\bibitem{hendrycks2021mmlu}
Dan Hendrycks, Collin Burns, Steven Basart, Andy Zou, Mantas Mazeika, Dawn Song,
  and Jacob Steinhardt, \emph{Measuring massive multitask language
  understanding}, International Conference on Learning Representations, 2021.

\bibitem{howard2020supermartingale}
Steven~R. Howard, Aaditya Ramdas, Jon McAuliffe, and Jasjeet Sekhon,
  \emph{Time-uniform Chernoff bounds via nonnegative supermartingales},
  Probability Surveys \textbf{17} (2020), 257--317.

\bibitem{howard2021confseq}
\bysame, \emph{Time-uniform, nonparametric, nonasymptotic confidence
  sequences}, Annals of Statistics \textbf{49} (2021), no.~2, 1055--1080.

\bibitem{knott1984optimal}
Martin Knott and Cyril~S. Smith, \emph{On the optimal mapping of distributions},
  Journal of Optimization Theory and Applications \textbf{43} (1984), 39--49.

\bibitem{lee2024cast}
Bruce~W. Lee, Inkit Padhi, Karthikeyan~Natesan Ramamurthy, Erik Miehling,
  Pierre Dognin, Manish Nagireddy, and Amit Dhurandhar, \emph{Programming
  refusal with conditional activation steering}, International Conference on
  Learning Representations, 2025.

\bibitem{li2023iti}
Kenneth Li, Oam Patel, Fernanda Vi\'egas, Hanspeter Pfister, and Martin
  Wattenberg, \emph{Inference-time intervention: Eliciting truthful answers from
  a language model}, Advances in Neural Information Processing Systems 36, 2023.

\bibitem{liu2026probe}
Ke Li, Chongzhe Zhang, Zifan Zeng, Feng Liu, Qunli Zhang, and Zheng Hu,
  \emph{Calibrating overconfidence without sacrificing confidence:
  Probe-conditioned head intervention for LLMs}, arXiv:2606.09876, 2026.

\bibitem{luo2026yopo}
Ziyang Luo, Zhongyao Chu, Xinjie He, Youting Wang, Xukui Qin, Runxiong Wu, and
  Yan-Syuan Chen, \emph{You only pass once: answering and abstaining together
  in a single forward pass of a frozen language model}, arXiv:2608.14465, 2026.

\bibitem{mishra2026nonsurjective}
Aayush Mishra, Daniel Khashabi, and Anqi Liu, \emph{Steered LLM activations are
  non-surjective}, arXiv:2604.09839, 2026.

\bibitem{nguyen2026alqr}
Julian Skifstad, Xinyue~Annie Yang, and Glen Chou, \emph{Local Linearity of
  LLMs Enables Activation Steering via Model-Based Linear Optimal Control},
  arXiv:2604.19018, 2026.

\bibitem{nguyen2026coast}
Tam Nguyen, Tu~Anh Nguyen, Sina Alemohammad, and Richard~G. Baraniuk,
  \emph{Minimizing collateral damage in activation steering}, arXiv:2605.01167,
  2026.

\bibitem{rimsky2024caa}
Nina Panickssery, Nick Gabrieli, Julian Schulz, Meg Tong, Evan Hubinger, and
  Alexander~Matt Turner, \emph{Steering Llama 2 via contrastive activation
  addition}, Proceedings of the 62nd Annual Meeting of the Association for
  Computational Linguistics, 2024.

\bibitem{rodriguez2025act}
Pau Rodriguez, Arno Blaas, Michal Klein, Luca Zappella, Nicholas Apostoloff,
  Marco Cuturi, and Xavier Suau, \emph{Controlling language and diffusion models
  by transporting activations}, International Conference on Learning
  Representations, 2025.

\bibitem{rodriguez2025lineas}
Pau Rodriguez, Michal Klein, Eleonora Gualdoni, Arno Blaas, Luca Zappella, Marco
  Cuturi, and Xavier Suau, \emph{LinEAS: End-to-end learning of activation
  steering with a distributional loss}, arXiv:2503.10679, 2025.

\bibitem{scalena2024dynamic}
Daniel Scalena, Gabriele Sarti, and Malvina Nissim, \emph{Multi-property
  steering of language models with dynamic activation composition}, Proceedings
  of BlackboxNLP, 2024.

\bibitem{singh2024mimic}
Shashwat Singh, Shauli Ravfogel, Jonathan Herzig, Roee Aharoni, Ryan Cotterell,
  and Ponnurangam Kumaraguru, \emph{Representation surgery: Theory and practice
  of affine steering}, Proceedings of the 41st International Conference on
  Machine Learning, 2024.

\bibitem{steerable2026}
\emph{When is your LLM steerable? Predicting steering success from internal
  states}, arXiv:2606.11599, 2026.

\bibitem{tan2024generalisation}
D.~Tan, D.~Chanin, A.~Lynch, B.~Paige, D.~Kanoulas, A.~Garriga-Alonso, R.~Kirk.
  \emph{Analysing the generalisation and reliability of steering vectors}.
  NeurIPS, 2024.

\bibitem{braun2025unreliability}
J.~Braun, C.~Eickhoff, D.~Krueger, S.~A.~Bahrainian, D.~Krasheninnikov.
  \emph{Understanding (un)reliability of steering vectors in language models}.
  ICLR Workshop on Foundation Models in the Wild, 2025.

\bibitem{templeton2024scaling}
Adly Templeton, Tom Conerly, Jonathan Marcus, et al., \emph{Scaling
  monosemanticity: Extracting interpretable features from Claude 3 Sonnet},
  Transformer Circuits Thread, 2024.

\bibitem{turner2023actadd}
Alexander~Matt Turner, Lisa Thiergart, Gavin Leech, David Udell, Juan~J.
  Vazquez, Ulisse Mini, and Monte MacDiarmid, \emph{Steering language models
  with activation engineering}, arXiv:2308.10248, 2023.

\bibitem{vogels2025ids}
Arthur Vogels, Benjamin Wong, Yann Choho, Annabelle Blangero, and Milan Bhan,
  \emph{In-distribution steering: Balancing control and coherence in language
  model generation}, arXiv:2510.13285, 2025.

\bibitem{vllmlens}
UK AI Safety Institute, \emph{vllm-lens: activation extraction and steering for
  vLLM}, \url{https://github.com/UKGovernmentBEIS/vllm-lens}, 2026.

\bibitem{wang2025sadi}
Weixuan Wang, Jingyuan Yang, and Wei Peng, \emph{Semantics-adaptive activation
  intervention for LLMs via dynamic steering vectors}, International Conference
  on Learning Representations, 2025.

\bibitem{wu2025axbench}
Zhengxuan Wu, Aryaman Arora, Atticus Geiger, Zheng Wang, Jing Huang, Dan
  Jurafsky, Christopher~D. Manning, and Christopher Potts, \emph{AxBench:
  Steering LLMs? Even simple baselines outperform sparse autoencoders},
  Proceedings of the 42nd International Conference on Machine Learning, 2025.

\bibitem{zhang2026wired}
Tianyi Zhao, Yinhan He, Wendy Zheng, Yujie Zhang, and Chen Chen, \emph{Wired
  for overconfidence: A mechanistic perspective on inflated verbalized
  confidence in LLMs}, Conference on Language Modeling (COLM), 2026,
  arXiv:2604.01457.

\bibitem{zou2023repe}
Andy Zou, Long Phan, Sarah Chen, James Campbell, et al., \emph{Representation
  engineering: A top-down approach to AI transparency}, arXiv:2310.01405, 2023.

\end{thebibliography}

\appendix

\section{Post-review stress-test boundary}
\label{app:postreview}

The confirmatory program preserved every planned cell in the denominator.
Mistral/TruthfulQA improved calibrated MC2 by $+.153$ with adjusted 95\% CI
$[+.108,+.197]$, but its MC1 interval $[-.059,+.122]$ crossed the
predeclared joint gate; Granite/TruthfulQA stopped at observer eligibility.
Five of six BBQ/ETHICS observer cells did not advance. Mistral/ETHICS did
advance and improved pooled calibrated MC2 by $+.107$ $[+.090,+.124]$, but
its label-1 MC1 changed by $-.093$ $[-.124,-.071]$. Thus the OLMo-2 result is
a confirmed architecture transfer, while model-wide and behavior-wide
generalization remain unsupported. All intervals use 10,000 paired BCa
resamples with the predeclared Bonferroni family.

\section{Full proofs}
\label{app:proofs}

\textbf{Theorem~\ref{thm:ot}.} By locality, $T(h)-h\in\mathrm{span}(V)$
$\mu$-a.s., hence $\lVert T(h)-h\rVert^2=\lVert VT(h)-Vh\rVert^2$ and
$C(T)=\mathbb{E}\lVert VT(h)-Vh\rVert^2$. Let $\gamma$ be the joint law of
$(Vh,\,VT(h))$; by the goal constraint $\gamma$ is a coupling of $(\mu_V,\nu)$
and $C(T)=\int\lVert y-x\rVert^2\,d\gamma\ge\Wtwo^2(\mu_V,\nu)$. Since
$\mu_V$ is absolutely continuous with respect to Lebesgue measure on
$\mathbb{R}^k$ with finite second moment (for $k{=}1$ atomlessness already
suffices), Brenier's theorem gives a unique optimal coupling induced by a map
$S=\nabla\varphi$, $\varphi$ convex; $T^*(h)=h+V^\top\!\big(S(Vh)-Vh\big)$ is admissible
and attains the bound. For any optimal $T$, the induced coupling must equal the
unique optimal one, i.e.\ $VT(h)=S(Vh)$ a.s. Since $S$ depends only on $Vh$, the
optimal intervention is $\sigma(Vh)$-measurable: access to $h_\perp$ is useless.
\qed

\textbf{Corollary~\ref{cor:additive}.} If $\nu=\mu_V*\delta_\alpha$ the monotone
map is $x\mapsto x+\alpha$. Conversely if the optimal map is $x+\alpha$ then
$\nu$ is the shift of $\mu_V$. For Gaussians with $\sigma_t\neq\sigma_s$, after
the best shift the remaining distance is
$\Wtwo^2(\mu_V*\delta_\alpha,\nu)=(\sigma_s-\sigma_t)^2>0$. Selectivity: let two
subpopulations have the same law of $p$. Any transfer function $g$ then induces
the same law of displacements $\Delta=(g(p)-p)\hat v$ on both, so no choice of
$g$ can address one and spare the other: the operator carries no information
that distinguishes them. If in addition the state-change probability depends on
$h$ only through $(p,\Delta)$ --- equivalently $h_\perp$ is conditionally
independent of the subpopulation given $p$ --- the two marginal change rates
coincide and selectivity is zero in expectation. Without that conditional
independence a residual gap can survive through $h_\perp$, but it is a property
of the downstream geometry and is not reachable by tuning $g$. Empirically it
is $+0.000\pm0.000$ at the dose that removes the behavior
(Table~\ref{tab:defaults}). \qed

\textbf{Theorem~\ref{thm:gate}.} (a) is the Neyman--Pearson lemma applied to the
class-conditional laws of $s$; monotone-likelihood-ratio families reduce the LR
test to thresholding, and for $\mathcal{N}(\mu_B,\Sigma)$ vs
$\mathcal{N}(\mu_G,\Sigma)$ the log-LR is an affine function of the trace
feature $z$, $\log\mathrm{LR}(z)=w^\top z+b$ with $w=\Sigma^{-1}(\mu_B-\mu_G)$;
thresholding the scalar $s(z)=w^\top z+b$ is therefore the LR test (LDA). (b) The concavified ROC contains the segment
path $(0,0)\to(\mathrm{FPR}^*,\mathrm{TPR}^*)\to(1,1)$ through the
Youden-optimal point; its area is
$\tfrac12(1+\mathrm{TPR}^*-\mathrm{FPR}^*)=\tfrac{1+J}{2}\le\mathrm{AUC}$. \qed

\textbf{Theorem~\ref{thm:joint}.} Decompose
$C(T)=P(\Gamma)\mathbb{E}[\lVert\Delta\rVert^2\mid\Gamma]+P(\Gamma^c)\mathbb{E}[\lVert\Delta\rVert^2\mid\Gamma^c]$.
Identity off $\Gamma$ zeroes the second term and is forced by the constraint;
the first term with the conditional goal constraint is exactly the
Theorem~\ref{thm:ot} problem for $\mu(\cdot\mid\Gamma)$. The constraints and
objective separate over the partition $\{\Gamma,\Gamma^c\}$, so the piecewise
optimum is global. \qed

\textbf{Proposition~\ref{prop:dose}.} The displacement of $T_\lambda$ is
$\lambda(s)^2\lVert S(Vh)-Vh\rVert^2$, so its expectation is
$\mathbb{E}_s[\lambda(s)^2 d(s)]$ because $\lambda$ is $s$-measurable. The same
measurability makes the gain $\mathbb{E}_s[\lambda(s)(\beta_1\pi(s)-\beta_0(1-\pi(s)))]$.
The objective is therefore an integral over $s$ of
$\lambda(\beta_1\pi-\beta_0(1-\pi))-\lambda^2 d/2c$, strictly concave in
$\lambda$ for $d(s)>0$, and the constraint set $[0,1]$ is separable in $s$, so
the integral is maximised by maximising the integrand at each $s$. Setting the
derivative to zero gives $\lambda=c((\beta_0+\beta_1)\pi-\beta_0)/d$ and the
clip enforces $[0,1]$; factoring $\beta_0+\beta_1$ out of the bracket gives
$\pi-\tau$ with $\tau=\beta_0/(\beta_0+\beta_1)$. Monotonicity in $\pi$ and the
zero set are immediate. As $c/d\to\infty$ the unclipped value diverges in sign
according to $\pi-\tau$, so $\lambda^*\to\mathbf{1}[\pi>\tau]$ wherever
$\pi(s)\ne\tau$, which is $P$-a.s.\ when $\pi$ has no atom at $\tau$; at any
finite $c/d$ the objective at $\lambda^*$ strictly exceeds its value at the
indicator on the event $\{0<\lambda^*<1\}$, which then has positive
probability. \qed

\textbf{Theorem~\ref{thm:seq}.} (a) The first $t$ positions generate a
$\sigma$-algebra $\mathcal{F}_t$; a gate measurable with respect to
$\mathcal{F}_t$ is a test between the class-conditional laws of $z_t$, so
Theorem~\ref{thm:gate}(a) applies verbatim and its concavified-ROC bound gives
the ceiling $2\,\mathrm{AUC}(z_t)-1$. Nothing is assumed about how $z_t$
relates to $z_\infty$: the head is fitted on $z_t$ itself. (b) Write $X_r$ for
the gate contribution of position $r$ and $D_r=X_r-s_\infty$. By assumption
$(D_r)$ is a martingale difference sequence with conditionally
$\sigma$-sub-Gaussian increments, so
$\exp(\lambda\sum_{r\le t}D_r-\lambda^2\sigma^2t/2)$ is a nonnegative
supermartingale for every $\lambda$~\cite{howard2020supermartingale}. Stitching
these over geometric epochs and applying Ville's inequality gives the
finite-law-of-the-iterated-logarithm boundary and
$P\big(\exists t:\ |s_t-s_\infty|>B_t(\delta)\big)\le\delta$
(\cite{howard2021confseq}, Eq.~11 at $\eta{=}2$, $s{=}1.4$). We measure
$\sigma$ rather than assume it: with $\lVert w\rVert{=}1$ and calibration
activation norms bounded by $R$, Hoeffding's lemma makes each contribution
$R$-sub-Gaussian, and we fit the working $\sigma$ to the deviation profile of
the extraction traces. On $\{s_\infty\le\tau\}$ the firing event
$s_t>\tau+B_t(\delta)$ implies $s_t-s_\infty>B_t(\delta)$, which that bound
excludes except on a set of measure $\delta$; on the complement,
$s_t\ge\tau+\Delta-B_t(\delta)$, which exceeds $\tau+B_t(\delta)$ as soon as
$B_t(\delta)\le\Delta/2$, and solving
$\sigma\sqrt{2\log(\log 2t/\delta)/t}\le\Delta/2$ gives
$t=O(\sigma^2\Delta^{-2}\log\log(1/\delta))$. (c) The reader is registered on a
layer at or above the actor, so within one forward pass the actor for position
$t$ observes the state accumulated over positions $<t$; the running mean needs
one scalar sum and one counter per sequence, and the edit at position $t$
changes only the keys and values written at position $t$, which have not yet
been consumed. \qed

\textbf{Proposition~\ref{prop:metric}.} We argue in the latent space rather
than through a global change of variables, because a general $M$ does not
preserve the duality between the measurement matrix $V$ and the displacement
subspace $\mathrm{span}(V^\top)$. Locality writes $\Delta h=V^\top a$ with
$a\in\mathbb{R}^k$, and since $VV^\top=I_k$ the goal constraint fixes
$a=S(Vh)-Vh$ for whichever map $S$ is used. Hence
$\lVert\Delta h\rVert_M^2=a^\top(VMV^\top)a=\lVert S(q)-q\rVert_{M_V}^2$ with
$q=Vh$: minimising $C$ over admissible $T$ is exactly the Monge problem on
$\mathbb{R}^k$ with cost $\lVert\cdot\rVert_{M_V}^2$. That cost is a Euclidean
square in the coordinates $z=M_V^{1/2}q$ --- a change of variables internal to
$\mathbb{R}^k$, so the locality constraint is untouched --- and Brenier's theorem
in $z$ gives the unique optimal $\nabla\varphi$; mapping back yields
$S_M=M_V^{-1/2}\!\circ\nabla\varphi\circ M_V^{1/2}$. For $M=I$, $M_V=I_k$ and
$S_M$ is Theorem~\ref{thm:ot}'s map. \qed

\textbf{Proposition~\ref{prop:compute}.} \emph{Lower bound:} let $\hat v_j$ be
dense ($(\hat v_j)_i\neq0$ for all $i$) and $g$ non-constant in $p_j$. If an
algorithm leaves coordinate $i$ unread, an adversary perturbs $h_i$; this
changes $p_j=\langle h,\hat v_j\rangle$ and hence, on a set of positive
measure, the correct output $T(h)$, while the algorithm's output is unchanged;
contradiction. Learned diff-in-means directions are dense almost surely.
\emph{Upper bound:} PTS computes $p=Vh$ ($kd$ multiply-adds), a monotone
$K$-knot lookup ($O(\log K)$ by binary search), and $h+V^\top(g(p)-p)$ ($kd$
more): $\Theta(kd+\log K)$ total. Memory $\Theta(kd+kK)$: reproducing a target
marginal at quantile resolution $1/K$ requires $\Omega(K)$ parameters, and
specifying the subspace requires $kd$. \qed

\section{GSM8K MCQ construction}
\label{app:gsm}
Each GSM8K test item with integer gold $g$ receives three deterministic
distractors drawn (seeded by item index) from
$\{g\pm1, g\pm2, g\pm10, 2g, \lfloor g/2\rfloor, \mathrm{round}(1.5g),
\mathrm{round}(0.8g), g\pm\mathrm{U}\{3..9\}\}$, deduplicated, shuffled with the
gold into four options. This converts free-form arithmetic into selection while
keeping the reasoning burden; we report transfer numbers on it and do not claim
GSM8K free-form results.

\section{The behavior across domains and scales (measurement study)}
\label{sec:measure}

\textbf{Overconfidence is model $\times$ instrument (2B).} Five confidence
readouts on identical MMLU questions ($n{=}1000\times3$ seeds) disagree
categorically: ECE ranges from $.391\,[.374,.409]$ (letter-logit M2, the
literature default) through $.249$/$.177$ (self-report/answer-distribution)
to $.104\,[.086,.120]$ (per-option yes/no M4); disjoint CIs; the
instruments even disagree on the answer (all five agree on $28.9\%$ of
questions). Hence our protocol requires every steering claim to hold under
multiple readouts. \textbf{Domains:} on MMLU the model is overconfident
(52/300 \ocw{}), on GSM8K-MC and GPQA \emph{under}confident (6 and 27
\ocw{}); the failure mode is domain-dependent, and the input-dependent
operators inherit exactly this dependence (Table~\ref{tab:transfer}).

\textbf{Scale replication (9B).} The full pipeline was replicated on
\texttt{gemma-2-9b-it} (42 layers, $d{=}3584$, layer 24) over the same
domains. Method-dependence survives scale (MMLU ECE $.126$--$.266$, GPQA
$.177$--$.499$, same instrument ranking); scale \emph{polarizes} rather than
cures (GSM8K nearly solved and calibrated, gap $-0.015$; GPQA confidence
grows faster than knowledge, M2 conf $.86$ at acc $.36$); the confidence
cone replicates (cosines $0.63$--$0.83$); confidence becomes \emph{more}
linearly readable (detector AUROC to $0.83$ vs $0.62$--$0.67$ at 2B), so the
Thm.~\ref{thm:gate}b ceiling rises with scale; and steering remains
one-directional, as Remark~1 predicts.


\section{What each method does: measured movement of real traces}
\label{app:plane}

\begin{figure}[!t]
\centering
\begin{tikzpicture}
\matrix[column sep=3mm, row sep=2mm] {
\begin{axis}[qpanel, title={additive steering (prior standard)},
  ylabel={\textbf{wrong \& overconfident}\\[-1pt]\scriptsize detection projection $\langle h,\hat u\rangle$},
  ylabel style={cadd!80!black, align=center}]
\qdecor
\addplot[arocw] table[col sep=comma]{pic/viz/cls_additive_ocw_ar.csv};
\node[anchor=south west, font=\tiny, align=left, fill=white, fill opacity=.9, text opacity=1, inner sep=1.5pt]
  at (axis cs:-33,-52) {moves all 50/50, far past the\\calibrated zone (off-manifold)};
\end{axis} &
\begin{axis}[qpanel, title={gated ablation (ours)}]
\qdecor
\addplot[arocw] table[col sep=comma]{pic/viz/cls_ours_gatedabl_ocw_ar.csv};
\addplot[stillocw] table[x=x,y=y,col sep=comma]{pic/viz/cls_ours_gatedabl_ocw_still.csv};
\node[anchor=south west, font=\tiny, align=left, fill=white, fill opacity=.9, text opacity=1, inner sep=1.5pt]
  at (axis cs:-33,-52) {44/52 move into the calibrated\\zone; 8 below the gate stay (dots)};
\end{axis} \\
\begin{axis}[qpanel, ylabel={\textbf{right \& confident}\\[-1pt]\scriptsize detection projection $\langle h,\hat u\rangle$},
  ylabel style={cclamp!80!black, align=center},
  xlabel={confidence projection $\langle h,\hat v\rangle$}]
\qdecor
\addplot[arcr] table[col sep=comma]{pic/viz/cls_additive_cr_ar.csv};
\node[anchor=south west, font=\tiny, align=left, fill=white, fill opacity=.9, text opacity=1, inner sep=1.5pt]
  at (axis cs:-33,-52) {collateral damage: all 76/76\\correct answers dragged out too};
\end{axis} &
\begin{axis}[qpanel, xlabel={confidence projection $\langle h,\hat v\rangle$}]
\qdecor
\addplot[arcr] table[col sep=comma]{pic/viz/cls_ours_gatedabl_cr_ar.csv};
\addplot[stillcr] table[x=x,y=y,col sep=comma]{pic/viz/cls_ours_gatedabl_cr_still.csv};
\node[anchor=south west, font=\tiny, align=left, fill=white, fill opacity=.9, text opacity=1, inner sep=1.5pt]
  at (axis cs:-33,-52) {30 untouched (dots); the 46 gated\\barely move, 54\% keep their state};
\end{axis} \\
};
\end{tikzpicture}
\caption{Each dot is one reasoning trace before steering, an arrow is its
measured movement after steering on real rollouts, and a bold dot without an
arrow was left untouched. The horizontal axis is the confidence projection
$\langle h,\hat v\rangle$ and the vertical axis the detection projection
$\langle h,\hat u\rangle$. The pink band is gate ON and the green ellipse is
the $2\sigma$ contour of the calibrated extraction-split states
(Appendix~\ref{app:defs}). Additive steering (left) overshoots all 50 \ocw{} far past the
calibrated zone and drags all 76 \crok{} along. Gated ablation (right) moves
\ocw{} into the zone and leaves 30 \crok{} bitwise untouched.}
\label{fig:plane}
\end{figure}

\subsection*{Where the benchmarks live, and how we compare with the label oracle}
Figure~\ref{fig:benchgeo} (left) plots each transfer benchmark's baseline
traces in the same plane: MMLU and ARC carry a real overconfident-wrong mass
in the gated (upper-right) region, while GSM8K and GPQA have almost none.
That is the geometric reason Table~\ref{tab:transfer} shows activity on the
former and near-identity on the latter. Figure~\ref{fig:benchgeo} (right)
compares our intervention to the \emph{label oracle}, the conditional
steering that moves exactly the overconfident-wrong points onto the calibrated
centroid: our gate catches $41/50$ of the oracle's targets
(recall $0.82$), misses $9$, and false-fires on $46/76$ confident-right
points (which the tail-local action then mostly spares: $25/46$ keep their
state). The remaining gap to the oracle is exactly the detection gap of
Thm.~\ref{thm:gate}b.

\begin{figure}[!t]
\centering
\begin{tikzpicture}
\matrix {
\begin{axis}[planepanel, width=.5\linewidth, title={ARC-Challenge},
  ylabel={proj.\ on $\hat u$}]
\addplot[ptrest] table[x=x,y=y,col sep=comma]{pic/viz/pl_arc_ncw.csv};
\addplot[ptrest] table[x=x,y=y,col sep=comma]{pic/viz/pl_arc_ncr.csv};
\addplot[ptcr] table[x=x,y=y,col sep=comma]{pic/viz/pl_arc_cr.csv};
\addplot[ptocw] table[x=x,y=y,col sep=comma]{pic/viz/pl_arc_ocw.csv};
\addplot[gray, dashed, thick] coordinates {(-35,-10.08) (95,-10.08)};
\end{axis} &
\begin{axis}[planepanel, width=.5\linewidth, title={GSM8K-MC}]
\addplot[ptrest] table[x=x,y=y,col sep=comma]{pic/viz/pl_gsm8k_ncw.csv};
\addplot[ptrest] table[x=x,y=y,col sep=comma]{pic/viz/pl_gsm8k_ncr.csv};
\addplot[ptcr] table[x=x,y=y,col sep=comma]{pic/viz/pl_gsm8k_cr.csv};
\addplot[ptocw] table[x=x,y=y,col sep=comma]{pic/viz/pl_gsm8k_ocw.csv};
\addplot[gray, dashed, thick] coordinates {(-35,-10.08) (95,-10.08)};
\end{axis} \\
};
\end{tikzpicture}\\[4pt]
\begin{tikzpicture}
\begin{axis}[planepanel, width=.8\linewidth, height=5.6cm,
  title={ours vs the label oracle (MMLU)},
  xlabel={proj.\ on $\hat v$ (confidence)}, ylabel={proj.\ on $\hat u$ (detect)}]
\addplot[ptrest] table[x=x,y=y,col sep=comma]{pic/viz/pl_mmlu_ncw.csv};
\addplot[ptrest] table[x=x,y=y,col sep=comma]{pic/viz/pl_mmlu_ncr.csv};
\addplot[ptcr] table[x=x,y=y,col sep=comma]{pic/viz/pl_mmlu_cr.csv};
\addplot[-stealth, gray!60, dashed, quiver={u=\thisrow{u}, v=\thisrow{v}}]
  table[col sep=comma]{pic/viz/orc_oracle.csv};
\addplot[disp] table[col sep=comma]{pic/viz/orc_hits.csv};
\addplot[-stealth, gray!85!black, thin, quiver={u=\thisrow{u}, v=\thisrow{v}}]
  table[col sep=comma]{pic/viz/orc_ff.csv};
\addplot[only marks, mark=o, mark size=2.4pt, cadd, very thick]
  table[x=x,y=y,col sep=comma]{pic/viz/orc_miss.csv};
\end{axis}
\end{tikzpicture}
\caption{Top: transfer benchmarks in the behavioral plane (gate threshold
dashed): ARC has an overconfident-wrong tail above the gate, GSM8K almost
none, and transfer activity follows the geometry. Bottom: gray dashed arrows
$=$ the label oracle (move every overconfident-wrong point to the calibrated
centroid); pink $=$ our measured moves on the same points ($45/52$ caught);
hollow orange circles $=$ the $7$ misses; dark-gray thin arrows $=$ false fires on
confident-right (mostly state-preserving). The oracle--ours difference is the
detection gap that bounds selectivity (Thm.~\ref{thm:gate}b).}
\label{fig:benchgeo}
\end{figure}

\section{Sample traces before and after gated ablation}
\label{app:samples}

On the 223/300 gated questions, gated ablation does not induce hedging:
hedge markers \emph{drop} from $0.25$ to $0.13$ per trace and length from
$86$ to $71$ words, $29$ of the $50$ \ocw{} leave that state and $12$ come
back correct. Typical
outcomes: an associative wrong argument replaced by an explicit
derivation, better-calibrated reasoning with the same letter, and a flagged
\crok{} rephrased but preserved (strong additive or CAST doses instead yield
repetition loops). The excerpts below are truncated, and complete rollouts
ship with the released artifact.\looseness=-1

\begin{table}[!t]\footnotesize\setlength{\tabcolsep}{4pt}
\centering
\resizebox{\linewidth}{!}{%
\begin{tabular}{l|ccccc|cc}
\hline
condition (MMLU, M2 traces) & words & closes & boxed & distinct-3 & degen.\% & hedge & certainty \\
\hline
unsteered & 86 & 0.98 & 0.91 & 0.91 & 2.7 & 0.25 & 0.12 \\
additive $-0.75$ \cite{rimsky2024caa} & 186 & 0.28 & 0.20 & 0.31 & 76.7 & 0.08 & 0.29 \\
clamp $q_{50}$ (ours) & 76 & 0.99 & 0.94 & 0.91 & 2.7 & 0.18 & 0.10 \\
gated ablation (ours) & 71 & 0.98 & 0.95 & 0.90 & 3.7 & 0.13 & 0.15 \\
sequential gate (ours) & 79 & 0.98 & 0.93 & 0.91 & 2.3 & 0.24 & 0.11 \\
\hline

\end{tabular}}
\caption{What each operator does to the reasoning trace. ``Closes'' is the
fraction of traces that emit \texttt{</think>} and ``boxed'' the fraction that
emit an answer; degen.\% counts traces with distinct-3 below $0.5$; hedge and
certainty are marker counts per trace from the fixed lexicons of
Appendix~\ref{app:defs}.}
\label{tab:traces}
\end{table}

Across the 226/300 gated questions, the intervention does \emph{not} inject
hedging vocabulary --- hedge-marker frequency \emph{drops} ($0.57\to0.27$ per
trace) --- and traces get shorter ($95\to78$ words). What changes is the mode
of reasoning: assertive shortcut arguments are replaced by definitional,
evidence-first ones, and the final answer changes on 83/226 gated questions
(13 overconfident-wrong become \emph{correct} under M4; 32/48 gated
confident-right survive intact). Ablating the confidence projection during
the \texttt{think} phase appears to remove a ``commit early'' pressure rather
than adding doubt. Three representative pairs (traces truncated):

\noindent\textbf{Overconfident-wrong $\to$ corrected} (\texttt{mmlu-4592}; gold C
, M4 state overconfident wrong $\to$ nonconfident wrong). The baseline trace jumps from a vague association to an answer; the steered trace states the definition and derives the (gold) answer.
\begin{quote}\scriptsize\emph{Question:} If the price of firm A's cell phone service rises by 5 percent and the quantity demanded for firm B's cell phone service increases by 10 percent, we can say that Options: A. demand for firm B is price elastic. B. supply\\[2pt]
\textbf{before:} \texttt{\textless{}think\textgreater{}The question asks about the relationship between the price of firm A's service and the quantity demanded for firm B's service. We know that the price of firm A's service increased and the quantity demanded for firm B's service increased. This suggests that the demand for firm B's service is likely to be price elastic. \textless{}/think\textgreater{} \textbackslash\{\}boxed\{A\}} $\to$ \texttt{A}\\[2pt]
\textbf{after:} \texttt{\textless{}think\textgreater{}The cross-price elasticity of demand measures the percentage change in the quantity demanded of one good in response to a percentage change in the price of another good. Since the price of firm A's service increased and the quantity demanded for firm B increased, the cross-price elasticity is positive. We can conclude that the two goods are substitutes. \textless{}/think\textgreater{} \textbackslash\{\}boxed\{C\}} $\to$ \texttt{C}
\end{quote}

\noindent\textbf{Overconfident-wrong $\to$ calibrated-right} (\texttt{mmlu-12923}; gold C
, M4 state overconfident wrong $\to$ nonconfident right). The baseline burns its budget re-asserting; the steered trace weighs all four options tersely and lands on the same letter with calibrated confidence.
\begin{quote}\scriptsize\emph{Question:} The best conclusion that can be drawn about the four-day ( compressed) workweek is that it has: Options: A. long-lasting positive effects on attitudes and productivity. B. long-lasting negative effects on attitudes and p\\[2pt]
\textbf{before:} \texttt{\textless{}think\textgreater{}The question asks for the best conclusion about the four-day workweek. We need to consider the potential effects on attitudes and productivity. There's no mention of specific evidence or studies, so we need to consider the general impact of a compressed workweek. \textless{}/think\textgreater{} \textless{}think\textgreater{}Option A is too strong. We don't have enough information to say it has long-lasting positive effects on attitudes and productivity. \textless{}/think\textgreater{} \textless{}think\textgreater{}Option B is too strong. We don't have enough information to say it has long-las} $\to$ \texttt{D}\\[2pt]
\textbf{after:} \texttt{\textless{}think\textgreater{}The question asks for the best conclusion about the four-day workweek. We need to consider the potential effects on attitudes and productivity. Option A is too strong, as we don't have enough information to guarantee long-lasting positive effects. Option B is too strong, as we don't have enough information to guarantee long-lasting negative effects. Option C is a possibility, but we don't know if it's the best conclusion. Option D is a possibility, but we don't know if it's the best conclusion. \textless{}/think\textgreater{} \textbackslash\{\}} $\to$ \texttt{D}
\end{quote}

\noindent\textbf{Confident-right, gated but spared} (\texttt{mmlu-5125}; gold B
, M4 state confident right $\to$ confident right). A false-flagged correct answer: the steered reasoning changes wording, the answer improves/stays correct, and the confident-right state survives.
\begin{quote}\scriptsize\emph{Question:} Which one of the following best states Ivan Pavlov's accidental discovery? Options: A. While experimenting with rats, Pavlov discovered that if he repeatedly placed pellets of food in one side of a "T-maze," the rats wou\\[2pt]
\textbf{before:} \texttt{\textless{}think\textgreater{}The question asks for Pavlov's accidental discovery. Let's look at the options and see which one describes an accidental discovery. Option B describes a learned response to a stimulus, which is not an accidental discovery. Option C describes a cat learning to press a lever, which is a learned behavior. Option D describes a pigeon learning to associate a key with food. Option A describes a rat learning to run to a specific side of a maze. This is a classic example of classical conditioning. Therefore,} $\to$ \texttt{A}\\[2pt]
\textbf{after:} \texttt{\textless{}think\textgreater{}The question asks for the accidental discovery of Pavlov. Option B is the best choice because it describes Pavlov's accidental discovery of the dog's salivation response to the sight and sound of the attendant who brought them food. \textless{}/think\textgreater{} \textbackslash\{\}boxed\{B\}} $\to$ \texttt{B}
\end{quote}


\subsection*{One question, four methods}
The most instructive comparison holds the question fixed
(\texttt{mmlu-4592}, gold \texttt{C}; baseline answers \texttt{A},
overconfident-wrong). Our gated ablation produces the clean definitional trace
shown above and lands on \texttt{C}. The additive-action methods do something
qualitatively different --- their traces \emph{degenerate into verbatim
repetition loops}, the signature of support violation (Remark~6): the constant
shift pushes activations off the data manifold and generation cycles. Trace
length balloons from 351 to $\sim$1500 characters of looped text; this is not
deliberation but decoding pathology, invisible in aggregate confidence numbers
yet obvious in the rollouts. MiMiC (distribution-matched, unconditional) stays
fluent but keeps the wrong answer.

\noindent\textbf{additive $-0.75$ \cite{rimsky2024caa}} ($\to$ \texttt{A}):
\begin{quote}\scriptsize\ttfamily \textless{}think\textgreater{}The cross-price elasticity of the cell phone service is equal to the price of the cell phone service of the firm A and the quantity of the cell phone service of the firm B. The cross-price elasticity of the cell phone service is equal to the price of the cell phone service of the firm A and the quantity of the cell phone service of\end{quote}

\noindent\textbf{CAST-style gate$+$add \cite{lee2024cast}} ($\to$ \texttt{A}):
\begin{quote}\scriptsize\ttfamily \textless{}think\textgreater{}The cross-price elasticity of demand for firm B is (10/5) = 2. The cross-price elasticity of demand for firm A is (10/5) = 2. The cross-price elasticity of demand for firm A is (10/5) = 2. The cross-price elasticity of demand for firm B is (10/5) = 2. The cross-price elasticity of demand for firm A is (10/5) = 2. The cross-price el\end{quote}

\noindent\textbf{MiMiC \cite{singh2024mimic}} ($\to$ \texttt{A}):
\begin{quote}\scriptsize\ttfamily \textless{}think\textgreater{}The question asks about the relationship between the price of firm A's service and the quantity demanded for firm B's service. We know that the price of firm A's service increased and the quantity demanded for firm B's service increased. This suggests that the demand for firm B's service is\end{quote}

\noindent\textbf{gated ablation (ours)} ($\to$ \texttt{C}):
\begin{quote}\scriptsize\ttfamily \textless{}think\textgreater{}The cross-price elasticity of demand measures the percentage change in the quantity demanded of one good in response to a percentage change in the price of another good. Since the price of firm A's service increased and the quantity demanded for firm B increased, the cross-price elasticity is positive. We can conclude that the tw\end{quote}


\section{Where to read and where to act: the layer sweep}
\label{app:layers}

Table~\ref{tab:layers} refits both axes and the gate at nine depths of
\texttt{gemma-2-2b-it} on the same 400 extraction traces. Two quantities move
in opposite directions. Token-level confidence separation grows monotonically
with depth (Cohen's $d$ from $0.53$ to $0.70$), because late layers carry the
answer itself. Detector quality, the quantity Theorem~\ref{thm:gate}b turns
into a selectivity ceiling, peaks in the middle and collapses in the last
layers, where the residual stream has already committed. Layer~14 is inside
the plateau and layer~16 is its maximum; the gap between them,
$2\,\mathrm{AUROC}-1$ of $0.538$ against $0.599$, is why the sequential gate
reads two layers above the layer it edits (Remark~\ref{rem:depth}).

\begin{table}[!t]\footnotesize\setlength{\tabcolsep}{5pt}
\centering
\begin{tabular}{l|cc|cc}
\hline
layer & Cohen's $d$ conf. & Cohen's $d$ \ocw{}/\crok{} & gate AUROC & ceiling $2\mathrm{AUC}{-}1$ \\
\hline
4 & +0.525 & -0.079 & 0.783 & +0.566 \\
7 & +0.535 & -0.093 & 0.756 & +0.511 \\
10 & +0.512 & -0.144 & 0.741 & +0.482 \\
12 & +0.516 & -0.150 & 0.742 & +0.485 \\
14 & +0.589 & -0.155 & 0.769 & +0.538 \\
16\,$\star$ & +0.631 & -0.141 & 0.799 & +0.599 \\
18 & +0.616 & -0.116 & 0.779 & +0.558 \\
22 & +0.675 & -0.110 & 0.791 & +0.582 \\
25 & +0.700 & -0.111 & 0.649 & +0.298 \\
\hline

\end{tabular}
\caption{Layer sweep on the extraction split ($400$ traces, both axes and the
LDA gate refit at each depth). $\star$ marks the depth the sequential gate
reads. The \ocw{}/\crok{} column is measured on the axis refit at that depth
and is therefore not the $d=0.049$ of Fig.~\ref{fig:cdf}, which is measured on
the layer-14 axis the operators act through.}
\label{tab:layers}
\end{table}

\subsection*{What a prefix budget costs the gate}

A head fitted on completed traces is a detector, and a detector is not
automatically a predictor: features that recognise a behavior in finished text
are known to forecast it poorly from a prefix~\cite{fpcg2026}. We therefore do
not reuse the whole-trace head online. For each budget $t$ we refit the LDA on
the mean projection over the first $t$ generated positions and measure its
AUROC on the same contrast (Table~\ref{tab:prefix}). The cost of deciding
early is smaller than the worry suggests: at layer~16 a head that sees eight
generated tokens reaches $0.736$ against $0.799$ for the whole trace, and by
128 tokens the gap is $0.016$. The dip at $t{=}32$ is sampling noise in a
two-feature fit on 400 traces and does not survive at either neighbour. Layer
16 dominates layer 14 at every budget, which is why the sequential gate reads
there.

\begin{table}[!t]\footnotesize\setlength{\tabcolsep}{6pt}
\centering
\begin{tabular}{l|cccccc}
\hline
gate AUROC (\ocw{} vs \crok{}) & $t{=}8$ & 16 & 32 & 64 & 128 & 256 \\
\hline
layer 14 & 0.708 & 0.712 & 0.666 & 0.707 & 0.747 & 0.766 \\
layer 16 & 0.736 & 0.740 & 0.691 & 0.738 & 0.783 & 0.795 \\
\hline

\end{tabular}
\caption{Detector quality as a function of the generated-token budget the gate
is allowed to see, with the head refitted for each budget. The whole-trace
values are $0.769$ at layer~14 and $0.799$ at layer~16.}
\label{tab:prefix}
\end{table}

\section{Equal-budget tuning}
\label{app:parity}

Every method in Table~\ref{tab:headtohead} receives at least as many
configurations as our own eight, on the same held-out MMLU tuning split of
$n{=}1200$ questions,
disjoint from the extraction split and from all five evaluation subsets, and we
select among them by one pre-declared rule: the highest selectivity among
configurations whose accuracy change is at least $-0.01$. The additive baseline
searches its dose, CAST searches threshold $\times$ dose, MiMiC searches layer
$\times$ shrinkage, Linear-AcT searches layer $\times\lambda$, our conditional
ablation searches the depth the gate reads $\times$ the quantile it fires at
$\times$ the local action, our transport variant searches the subspace
dimension $k$ $\times$ the gate quantile, and the sequential gate searches its
prefix budget $t^\star\times$ the gate quantile. We report each winner in
Table~\ref{tab:parity}, and those configurations, unchanged, produce the
evaluation-seed numbers in the main text. The split has to be that large. On a
$400$-question split six of the nine winners are different and the measured
selectivity of a fixed configuration moves by up to $0.131$
(Table~\ref{tab:split}), so the smaller split does not identify a
configuration, it samples one. Every search also runs at one batch size,
because batched greedy decoding is not invariant to it and the same MiMiC
configuration measures $+0.288$ at $64$ and $+0.235$ at $32$. Each method's search is paired
against a baseline generated in the same process, which is what config
selection needs; the searches ran on different GPUs and their baselines differ
by up to three \ocw{} out of 79, so we do not compare selectivities
\emph{across} methods on this split. That comparison is
Table~\ref{tab:headtohead}, where every method is paired against one shared
set of baseline rollouts per seed. We also report one configuration
that is not part of any selection, the NP gate composed with the full-space
affine map, because Theorem~\ref{thm:joint} predicts the gate composes with any
transport and that is the cheapest way to test the prediction on the strongest
unconditional action.

\begin{table}[!t]\footnotesize\setlength{\tabcolsep}{4pt}
\centering
\resizebox{\linewidth}{!}{%
\begin{tabular}{l|c|l|cccc}
\hline
method & budget & selected configuration & $\Delta$acc & ocw\_rm & cr\_keep & Sel \\
\hline
additive CAA \cite{rimsky2024caa} & 8 & \texttt{alpha-0.25} & +0.016 & 0.627 & 0.644 & +0.271 \\
CAST-style \cite{lee2024cast} & 8 & \texttt{tau\_cr\_q70\_alpha-0.375} & -0.005 & 0.428 & 0.756 & +0.184 \\
MiMiC \cite{singh2024mimic} & 9 & \texttt{L18\_reg0.1} & +0.030 & 0.468 & 0.795 & +0.263 \\
Linear-AcT \cite{rodriguez2025act} & 9 & \texttt{L14\_lam1.0} & +0.026 & 0.318 & 0.865 & +0.184 \\
gated local action (ours) & 8 & \texttt{gl14\_q30\_ablate} & +0.001 & 0.527 & 0.718 & +0.245 \\
gated $k$-dim transport (ours) & 8 & \texttt{k2\_q30} & +0.005 & 0.279 & 0.907 & +0.186 \\
score-proportional dose (ours) & 8 & \texttt{gl16\_q10\_kappa1.0} & +0.010 & 0.607 & 0.660 & +0.267 \\
sequential gate (ours, single pass) & 8 & \texttt{q50\_at32} & +0.014 & 0.219 & 0.917 & +0.136 \\
NP-gate $\times$ MiMiC (Thm.~\ref{thm:joint}) & 4 & \texttt{L14\_q30\_mimic} & +0.022 & 0.318 & 0.837 & +0.155 \\
\hline

\end{tabular}}
\caption{Equal-budget hyperparameter search on the tuning split.}
\label{tab:parity}
\end{table}

\begin{table}[!t]\footnotesize\setlength{\tabcolsep}{3pt}
\centering
\resizebox{\linewidth}{!}{%
\begin{tabular}{l|lc|lcc}
\hline
 & \multicolumn{2}{c|}{$n{=}400$ split} & \multicolumn{3}{c}{$n{=}1200$ split} \\
method & winner & Sel & winner & Sel & $\Delta$ \\
\hline
additive CAA \cite{rimsky2024caa} & \texttt{alpha-0.25} & +0.196 & \texttt{alpha-0.25} & +0.271 & +0.075 \\
CAST-style \cite{lee2024cast} & \texttt{tau\_cr\_q70\_alpha-0.75} & +0.182 & \texttt{tau\_cr\_q70\_alpha-0.375} & +0.184 & +0.003 \\
MiMiC \cite{singh2024mimic} & \texttt{L18\_reg0.01} & +0.288 & \texttt{L18\_reg0.1} & +0.263 & -0.025 \\
Linear-AcT \cite{rodriguez2025act} & \texttt{L10\_lam0.25} & +0.099 & \texttt{L14\_lam1.0} & +0.184 & +0.085 \\
gated ablation (tuned, ours) & \texttt{gl14\_q30\_ablate} & +0.232 & \texttt{gl14\_q30\_ablate} & +0.245 & +0.013 \\
gated $k$-dim transport (ours) & \texttt{k32\_q30} & +0.165 & \texttt{k2\_q30} & +0.186 & +0.021 \\
score-proportional dose (ours) & \texttt{gl16\_q10\_kappa0.5} & +0.265 & \texttt{gl16\_q10\_kappa1.0} & +0.267 & +0.003 \\
sequential gate (ours, single pass) & \texttt{q30\_at16} & +0.120 & \texttt{q50\_at32} & +0.136 & +0.015 \\
NP-gate $\times$ MiMiC (Thm.~\ref{thm:joint}) & \texttt{L18\_q30\_mimic} & +0.286 & \texttt{L14\_q30\_mimic} & +0.155 & -0.131 \\
\hline

\end{tabular}}
\caption{The same search on a tuning split three times larger. Six of the
nine winners change and the measured selectivity moves by up to $0.131$, so a
$400$-question split does not identify a configuration reliably. The paper
uses the $n{=}1200$ winners throughout.}
\label{tab:split}
\end{table}

\section{Design ablations of the operator family}
\label{app:ablation}

Every member of the family Theorem~\ref{thm:ot} defines was tuned under the
same budget and carried unchanged to the five evaluation subsets. The point of
the table is which design choices survive that carry. Raising the subspace
dimension does not: the tuned $k$ falls to $+0.065$ on the subsets from
$+0.186$ on the split. Making the action token-local rather than trace-local
does not either, $+0.158$ from $+0.257$, and neither does weighting the dose
by a posterior. Grading the dose by the gate margin does survive, $+0.231$
from $+0.267$, and so does the plain $k{=}2$ gated ablation, $+0.206$ from
$+0.245$, which loses the least of anything we tried. One member is worth reporting on its own terms. Theorem~\ref{thm:joint} gives
the cost-minimal gated intervention as the transport map of the law
\emph{conditioned} on the gate event, and every gated transport above uses the
marginal law instead. We fitted the conditional moments on the extraction
split, at each of the four gate quantiles, and compared the two maps under one
budget. The conditional map is not better: $+0.155$ against $+0.183$ at the
20th percentile, $+0.153$ against $+0.125$ at the 30th, and no systematic
ordering across the four. We read that as Theorem~\ref{thm:joint} delivering
what it claims and nothing else. It minimises displacement subject to reaching
the target law, which is not the same objective as maximising selectivity, and
selectivity here is bound by the detector rather than by the map
(Thm.~\ref{thm:gate}b). Five attempts to buy more removal by changing the
action --- a larger subspace, a full-rank map, a graded dose, a token-local
clamp and the conditional map --- each bought it at the same rate in
retention, which is what Cor.~\ref{cor:additive} predicts for a subpopulation
pair separated by $d=0.049$.

We report the members that did not carry because each one is an insight about
the selection problem rather than a catalogue entry. Together they say that
the more exposure or dose a configuration buys on a tuning split, the less of
it survives a change of questions, and that the design choice which survives
best is the one the theory singles out: a two-dimensional edit applied behind
a gate.

\begin{table}[!t]\footnotesize\setlength{\tabcolsep}{4pt}
\centering
\resizebox{\linewidth}{!}{%
\begin{tabular}{l|cccc}
\hline
member (M4, MMLU, seeds mean$\pm$std) & $\Delta$acc & ECE & cr\_keep & Sel \\
\hline
% seeds: 7,11,23,31,47 (n=300 each, M4 readout)
baseline & acc $0.473\pm0.026$ & $0.099\pm0.019$ & --- & --- \\
gated ablation (tuned), $k{=}2$ & $-0.015\pm0.036$ & $0.089\pm0.013$ & $0.713\pm0.050$ & $+0.206\pm0.059$ \\
gated $k$-dim transport & $-0.010\pm0.012$ & $0.100\pm0.011$ & $0.882\pm0.053$ & $+0.065\pm0.095$ \\
gated token-local clamp & $-0.008\pm0.053$ & $0.095\pm0.023$ & $0.721\pm0.050$ & $+0.158\pm0.085$ \\
score-proportional dose & $-0.005\pm0.054$ & $0.089\pm0.010$ & $0.650\pm0.056$ & $+0.231\pm0.120$ \\
posterior-weighted dose & $-0.003\pm0.006$ & $0.093\pm0.017$ & $0.969\pm0.027$ & $+0.042\pm0.050$ \\
gated full-rank transport & $+0.003\pm0.021$ & $0.108\pm0.022$ & $0.854\pm0.038$ & $+0.146\pm0.042$ \\
\hline

\end{tabular}}
\caption{Members of the projection-transport family under the equal-budget
protocol, each at the configuration it won on the $n{=}1200$ tuning split.}
\label{tab:ablation}
\end{table}

\section{Full result tables}
\label{app:fulltables}

\begin{table}[!t]\footnotesize\setlength{\tabcolsep}{4pt}
\centering
\resizebox{\linewidth}{!}{%
\begin{tabular}{l|cccc}
\hline
method (M4, MMLU, seeds mean$\pm$std) & $\Delta$acc & ECE & AURC & Sel \\
\hline
% seeds: 7,11,23,31,47 (n=300 each, M4 readout)
baseline & acc $0.473\pm0.026$ & $0.099\pm0.019$ & $0.404\pm0.029$ & --- \\
prompting & $-0.005\pm0.020$ & $0.118\pm0.018$ & $0.392\pm0.023$ & $+0.198\pm0.086$ \\
additive $-0.75$ \cite{rimsky2024caa} & $-0.087\pm0.053$ & $0.108\pm0.044$ & $0.556\pm0.056$ & $+0.000\pm0.000$ \\
CAST-style gate$+$add \cite{lee2024cast} & $-0.049\pm0.039$ & $0.099\pm0.021$ & $0.450\pm0.040$ & $+0.222\pm0.055$ \\
MiMiC full-space affine \cite{singh2024mimic} & $-0.001\pm0.041$ & $0.116\pm0.026$ & $0.401\pm0.048$ & $+0.200\pm0.107$ \\
Linear-AcT $\lambda{=}0.5$ \cite{rodriguez2025act} & $-0.006\pm0.036$ & $0.104\pm0.015$ & $0.412\pm0.018$ & $+0.120\pm0.080$ \\
Linear-AcT $\lambda{=}1.0$ \cite{rodriguez2025act} & $+0.001\pm0.032$ & $0.101\pm0.011$ & $0.410\pm0.027$ & $+0.146\pm0.141$ \\
clamp $q_{50}$ (ours) & $+0.001\pm0.034$ & $0.110\pm0.033$ & $0.408\pm0.020$ & $+0.145\pm0.039$ \\
quantile transport (ours) & $-0.001\pm0.029$ & $0.108\pm0.017$ & $0.408\pm0.015$ & $+0.119\pm0.086$ \\
gated clamp (ours) & $-0.001\pm0.039$ & $0.096\pm0.021$ & $0.404\pm0.024$ & $+0.224\pm0.065$ \\
gated ablation (ours) & $-0.003\pm0.050$ & $0.094\pm0.030$ & $0.395\pm0.038$ & $+0.279\pm0.090$ \\
LDA-gated ablation (ours) & $-0.009\pm0.032$ & $0.091\pm0.010$ & $0.402\pm0.028$ & $+0.198\pm0.094$ \\
online-gated ablation (ours, single pass) & $-0.000\pm0.016$ & $0.109\pm0.016$ & $0.405\pm0.040$ & $+0.035\pm0.048$ \\
\hline

\end{tabular}}
\caption{The same five evaluation subsets with every method at its default
configuration, that is before the equal-budget search of
Appendix~\ref{app:parity}.}
\label{tab:defaults}
\end{table}

\begin{table}[!t]\footnotesize\setlength{\tabcolsep}{3.4pt}
\centering
\resizebox{\linewidth}{!}{%
\begin{tabular}{l|cc|cc|c|c}
\hline
condition (M4 readout) & $\Delta$acc & 95\% CI & ECE & \ocw{} left & cr\_keep & Sel [95\% CI] \\
\hline
baseline & --- & --- & 0.076 & 50 & --- & --- \\
additive $-0.75$ \cite{rimsky2024caa} & $-0.150$ & $[-0.21,-0.08]$ & 0.052 & 0 & 0.000 & $+0.000$\,$[+0.00,+0.00]$ \\
ablation \cite{arditi2024refusal} & $+0.020$ & $[-0.04,+0.08]$ & 0.079 & 30 & 0.592 & $+0.232$\,$[+0.07,+0.40]$ \\
clamp $q_{50}$ & $+0.013$ & $[-0.04,+0.07]$ & 0.101 & 48 & 0.724 & $+0.124$\,$[-0.05,+0.30]$ \\
quantile transport & $-0.007$ & $[-0.04,+0.03]$ & 0.096 & 54 & 0.829 & $+0.029$\,$[-0.12,+0.18]$ \\
gated clamp $q_{50}$ & $+0.017$ & $[-0.03,+0.07]$ & 0.081 & 47 & 0.855 & $+0.195$\,$[+0.04,+0.35]$ \\
gated clamp $q_{70}$ & $-0.003$ & $[-0.04,+0.04]$ & 0.045 & 39 & 0.895 & $+0.175$\,$[+0.04,+0.32]$ \\
gated ablation & $+0.033$ & $[-0.02,+0.09]$ & 0.101 & 35 & 0.724 & $+0.304$\,$[+0.13,+0.47]$ \\
sequential-gate ablation & $+0.003$ & $[-0.03,+0.04]$ & 0.098 & 49 & 0.934 & $+0.074$\,$[-0.03,+0.19]$ \\
\hline

\end{tabular}}
\caption{MMLU, $n{=}300$, layer 14, M4 readout (the calibrated instrument).
The additive baseline removes all \ocw{} by destroying \crok{} (selectivity
$\approx0$, as Cor.~\ref{cor:additive} forces). Input-dependent transfer
functions achieve positive selectivity at an accuracy \emph{gain}. Same-sign
effects hold under the M2 readout (Appendix~\ref{app:fulltables}).}
\label{tab:main}
\end{table}


Full condition grids under both readouts. Columns: condition, readout,
$\Delta$accuracy with paired-bootstrap 95\% CI, ECE, count of \ocw{} after
steering, \crok{}-retention, selectivity with 95\% CI. \texttt{subbw} = ungated
subspace Bures--Wasserstein transport; \texttt{npbw} = LDA-gated subspace
transport (\texttt{self\_qXX} = gate threshold re-anchored, label-free, to the
target domain's own score quantile); \texttt{-think} = steering active only
during the \texttt{<think>} phase.

\begin{table}[!t]\scriptsize\setlength{\tabcolsep}{2.6pt}
\centering
\resizebox{\linewidth}{!}{%
\begin{tabular}{l|c|cc|cc|c}
\hline
condition & r/o & $\Delta$acc [CI] & ECE & \ocw{} & cr\_keep & Sel [CI] \\
\hline
baseline & m4 & --- & 0.076 & 50 & --- & --- \\
Linear-AcT $\lambda{=}0.5$ & m4 & $-0.007\,[-0.053,+0.037]$ & 0.084 & 50 & 0.83 & $+0.109$\,$[-0.04,+0.25]$ \\
Linear-AcT $\lambda{=}1.0$ & m4 & $+0.003\,[-0.037,+0.043]$ & 0.090 & 48 & 0.88 & $+0.182$\,$[+0.03,+0.33]$ \\
CAST-style gate$+$add & m4 & $-0.067\,[-0.130,-0.003]$ & 0.079 & 9 & 0.39 & $+0.215$\,$[+0.05,+0.37]$ \\
ablation & m4 & $+0.020\,[-0.040,+0.080]$ & 0.079 & 30 & 0.59 & $+0.232$\,$[+0.07,+0.40]$ \\
additive $-0.75$ & m4 & $-0.150\,[-0.213,-0.080]$ & 0.052 & 0 & 0.00 & $+0.000$\,$[+0.00,+0.00]$ \\
clamp $q_{30}$ & m4 & $+0.027\,[-0.030,+0.083]$ & 0.102 & 40 & 0.68 & $+0.184$\,$[+0.01,+0.36]$ \\
clamp $q_{50}$ & m4 & $+0.013\,[-0.043,+0.070]$ & 0.101 & 48 & 0.72 & $+0.124$\,$[-0.05,+0.30]$ \\
clamp $q_{70}$ & m4 & $+0.023\,[-0.030,+0.077]$ & 0.136 & 48 & 0.72 & $+0.144$\,$[-0.03,+0.31]$ \\
gated clamp $q_{50}$ & m4 & $+0.017\,[-0.033,+0.067]$ & 0.081 & 47 & 0.86 & $+0.195$\,$[+0.04,+0.35]$ \\
gated OT $q_{50}$ & m4 & $-0.023\,[-0.067,+0.020]$ & 0.095 & 50 & 0.84 & $+0.042$\,$[-0.10,+0.18]$ \\
gated clamp $q_{70}$ & m4 & $-0.003\,[-0.043,+0.040]$ & 0.045 & 39 & 0.89 & $+0.175$\,$[+0.04,+0.32]$ \\
gated OT $q_{70}$ & m4 & $-0.010\,[-0.037,+0.017]$ & 0.083 & 53 & 0.92 & $+0.001$\,$[-0.09,+0.10]$ \\
Gaussian-OT cal. & m4 & $+0.027\,[-0.010,+0.067]$ & 0.119 & 48 & 0.88 & $+0.062$\,$[-0.06,+0.19]$ \\
quantile-OT cal. & m4 & $+0.007\,[-0.030,+0.047]$ & 0.131 & 55 & 0.84 & $+0.002$\,$[-0.13,+0.14]$ \\
quantile-OT non-\ocw{} & m4 & $-0.007\,[-0.043,+0.030]$ & 0.096 & 54 & 0.83 & $+0.029$\,$[-0.12,+0.18]$ \\
MiMiC full-space affine & m4 & $+0.013\,[-0.037,+0.063]$ & 0.098 & 43 & 0.83 & $+0.249$\,$[+0.09,+0.40]$ \\
sequential-gate ablation & m4 & $+0.003\,[-0.027,+0.037]$ & 0.098 & 49 & 0.93 & $+0.074$\,$[-0.03,+0.19]$ \\
online\_d0p05\_q50\_L16\_t16\_ablate & m4 & $+0.017\,[-0.023,+0.057]$ & 0.124 & 53 & 0.79 & $-0.011$\,$[-0.15,+0.13]$ \\
online\_d0p05\_q50\_L16\_t64\_ablate & m4 & $-0.007\,[-0.027,+0.013]$ & 0.075 & 47 & 0.93 & $+0.054$\,$[-0.05,+0.16]$ \\
probegate-all\_q60\_ablate & m4 & $+0.013\,[-0.023,+0.053]$ & 0.105 & 42 & 0.74 & $+0.117$\,$[-0.04,+0.29]$ \\
probegate-all\_q60\_clamp\_q30 & m4 & $+0.027\,[-0.010,+0.063]$ & 0.109 & 46 & 0.78 & $+0.076$\,$[-0.08,+0.24]$ \\
probegate-all\_q60\_otq\_cal & m4 & $-0.013\,[-0.043,+0.017]$ & 0.088 & 48 & 0.92 & $+0.061$\,$[-0.05,+0.19]$ \\
probegate-all\_q75\_ablate & m4 & $+0.020\,[-0.013,+0.057]$ & 0.107 & 44 & 0.87 & $+0.188$\,$[+0.04,+0.34]$ \\
probegate-all\_q75\_clamp\_q30 & m4 & $+0.033\,[+0.000,+0.070]$ & 0.074 & 40 & 0.86 & $+0.195$\,$[+0.05,+0.35]$ \\
probegate-all\_q75\_otq\_cal & m4 & $-0.003\,[-0.030,+0.027]$ & 0.089 & 51 & 0.89 & $+0.035$\,$[-0.08,+0.16]$ \\
prompting & m4 & $-0.017\,[-0.073,+0.043]$ & 0.096 & 58 & 0.70 & $+0.097$\,$[-0.08,+0.26]$ \\
LDA-gated ablation & m4 & $+0.020\,[-0.013,+0.053]$ & 0.079 & 38 & 0.83 & $+0.209$\,$[+0.04,+0.37]$ \\
sweep\_lda-allq50\_clamp\_q30 & m4 & $+0.030\,[-0.003,+0.067]$ & 0.093 & 44 & 0.89 & $+0.175$\,$[+0.03,+0.33]$ \\
sweep\_lda-allq50\_clamp\_q50 & m4 & $+0.013\,[-0.020,+0.047]$ & 0.100 & 45 & 0.95 & $+0.187$\,$[+0.06,+0.33]$ \\
gated ablation & m4 & $+0.033\,[-0.020,+0.090]$ & 0.101 & 35 & 0.72 & $+0.304$\,$[+0.13,+0.47]$ \\
sweep\_ocwcr-crq50\_clamp\_q30 & m4 & $+0.007\,[-0.040,+0.053]$ & 0.078 & 48 & 0.84 & $+0.222$\,$[+0.07,+0.38]$ \\
gated clamp $q_{50}$ (sweep) & m4 & $+0.017\,[-0.033,+0.067]$ & 0.081 & 47 & 0.86 & $+0.195$\,$[+0.04,+0.35]$ \\
Linear-AcT (tuned) & m4 & $+0.003\,[-0.037,+0.043]$ & 0.090 & 48 & 0.88 & $+0.182$\,$[+0.03,+0.33]$ \\
additive CAA (tuned) & m4 & $+0.000\,[-0.057,+0.063]$ & 0.079 & 40 & 0.59 & $+0.232$\,$[+0.05,+0.41]$ \\
CAST-style (tuned) & m4 & $+0.007\,[-0.037,+0.053]$ & 0.076 & 33 & 0.76 & $+0.123$\,$[-0.04,+0.29]$ \\
gated full-rank transport (tuned, ours) & m4 & $+0.027\,[-0.010,+0.063]$ & 0.090 & 46 & 0.92 & $+0.181$\,$[+0.05,+0.32]$ \\
MiMiC (tuned) & m4 & $+0.013\,[-0.040,+0.070]$ & 0.110 & 56 & 0.75 & $+0.070$\,$[-0.09,+0.23]$ \\
sequential gate (tuned, ours) & m4 & $-0.007\,[-0.037,+0.023]$ & 0.114 & 49 & 0.88 & $+0.062$\,$[-0.06,+0.19]$ \\
gated local action (tuned, ours) & m4 & $+0.027\,[-0.020,+0.077]$ & 0.094 & 35 & 0.78 & $+0.276$\,$[+0.11,+0.45]$ \\
gated $k$-dim transport (tuned, ours) & m4 & $-0.027\,[-0.063,+0.013]$ & 0.095 & 57 & 0.82 & $-0.064$\,$[-0.19,+0.06]$ \\
tuned\_ourspost & m4 & $-0.003\,[-0.020,+0.013]$ & 0.075 & 48 & 0.99 & $+0.067$\,$[-0.01,+0.15]$ \\
score-proportional dose (tuned, ours) & m4 & $+0.043\,[-0.010,+0.100]$ & 0.091 & 33 & 0.66 & $+0.218$\,$[+0.04,+0.39]$ \\
gated token-local clamp (tuned, ours) & m4 & $+0.023\,[-0.023,+0.070]$ & 0.084 & 39 & 0.76 & $+0.263$\,$[+0.10,+0.44]$ \\
baseline & m2 & --- & 0.375 & 124 & --- & --- \\
Linear-AcT $\lambda{=}0.5$ & m2 & $+0.017\,[-0.013,+0.047]$ & 0.360 & 121 & 0.94 & $+0.054$\,$[-0.01,+0.12]$ \\
Linear-AcT $\lambda{=}1.0$ & m2 & $+0.007\,[-0.027,+0.040]$ & 0.374 & 124 & 0.92 & $+0.052$\,$[-0.02,+0.12]$ \\
CAST-style gate$+$add & m2 & $-0.110\,[-0.167,-0.050]$ & 0.227 & 89 & 0.60 & $+0.090$\,$[-0.03,+0.20]$ \\
ablation & m2 & $+0.013\,[-0.037,+0.063]$ & 0.344 & 120 & 0.84 & $+0.106$\,$[+0.01,+0.20]$ \\
additive $-0.75$ & m2 & $-0.193\,[-0.260,-0.127]$ & 0.219 & 98 & 0.38 & $-0.051$\,$[-0.16,+0.07]$ \\
clamp $q_{30}$ & m2 & $+0.010\,[-0.040,+0.057]$ & 0.354 & 122 & 0.85 & $+0.104$\,$[+0.01,+0.20]$ \\
clamp $q_{50}$ & m2 & $+0.007\,[-0.037,+0.050]$ & 0.368 & 125 & 0.88 & $+0.059$\,$[-0.02,+0.14]$ \\
clamp $q_{70}$ & m2 & $+0.003\,[-0.043,+0.047]$ & 0.380 & 126 & 0.87 & $+0.063$\,$[-0.02,+0.15]$ \\
gated clamp $q_{50}$ & m2 & $+0.023\,[-0.020,+0.067]$ & 0.355 & 119 & 0.90 & $+0.101$\,$[+0.02,+0.18]$ \\
gated OT $q_{50}$ & m2 & $+0.000\,[-0.027,+0.027]$ & 0.380 & 127 & 0.95 & $+0.009$\,$[-0.04,+0.06]$ \\
gated clamp $q_{70}$ & m2 & $+0.030\,[-0.007,+0.067]$ & 0.348 & 116 & 0.95 & $+0.116$\,$[+0.04,+0.19]$ \\
gated OT $q_{70}$ & m2 & $+0.010\,[-0.007,+0.027]$ & 0.370 & 121 & 0.99 & $+0.045$\,$[+0.00,+0.09]$ \\
Gaussian-OT cal. & m2 & $-0.003\,[-0.030,+0.023]$ & 0.380 & 128 & 0.95 & $-0.005$\,$[-0.05,+0.04]$ \\
quantile-OT cal. & m2 & $+0.010\,[-0.017,+0.040]$ & 0.368 & 125 & 0.95 & $+0.035$\,$[-0.02,+0.10]$ \\
quantile-OT non-\ocw{} & m2 & $+0.017\,[-0.010,+0.047]$ & 0.370 & 122 & 0.96 & $+0.055$\,$[-0.00,+0.12]$ \\
MiMiC full-space affine & m2 & $+0.020\,[-0.023,+0.063]$ & 0.360 & 120 & 0.88 & $+0.091$\,$[+0.01,+0.17]$ \\
sequential-gate ablation & m2 & $+0.007\,[-0.010,+0.027]$ & 0.363 & 123 & 0.98 & $+0.017$\,$[-0.02,+0.06]$ \\
sequential-gate clamp $q_{50}$ & m2 & $+0.000\,[-0.017,+0.017]$ & 0.377 & 124 & 0.98 & $+0.006$\,$[-0.03,+0.04]$ \\
online\_d0p05\_q50\_L16\_t16\_ablate & m2 & $+0.013\,[-0.010,+0.040]$ & 0.357 & 120 & 0.96 & $+0.053$\,$[-0.00,+0.11]$ \\
online\_d0p05\_q50\_L16\_t64\_ablate & m2 & $+0.010\,[-0.007,+0.027]$ & 0.369 & 121 & 0.99 & $+0.029$\,$[-0.01,+0.07]$ \\
probegate-all\_q60\_ablate & m2 & $+0.000\,[-0.033,+0.030]$ & 0.367 & 122 & 0.93 & $+0.048$\,$[-0.02,+0.12]$ \\
probegate-all\_q60\_clamp\_q30 & m2 & $-0.013\,[-0.047,+0.020]$ & 0.383 & 127 & 0.92 & $+0.020$\,$[-0.04,+0.09]$ \\
probegate-all\_q60\_otq\_cal & m2 & $+0.013\,[-0.007,+0.033]$ & 0.369 & 120 & 0.98 & $+0.047$\,$[+0.00,+0.10]$ \\
probegate-all\_q75\_ablate & m2 & $+0.000\,[-0.030,+0.030]$ & 0.365 & 123 & 0.93 & $+0.018$\,$[-0.04,+0.08]$ \\
probegate-all\_q75\_clamp\_q30 & m2 & $-0.007\,[-0.037,+0.023]$ & 0.378 & 126 & 0.92 & $+0.012$\,$[-0.05,+0.08]$ \\
probegate-all\_q75\_otq\_cal & m2 & $+0.007\,[-0.017,+0.030]$ & 0.367 & 123 & 0.97 & $+0.027$\,$[-0.02,+0.07]$ \\
prompting & m2 & $+0.020\,[-0.030,+0.073]$ & 0.344 & 119 & 0.85 & $+0.110$\,$[+0.01,+0.20]$ \\
LDA-gated ablation & m2 & $-0.007\,[-0.033,+0.020]$ & 0.365 & 125 & 0.94 & $+0.005$\,$[-0.05,+0.06]$ \\
sweep\_lda-allq50\_clamp\_q30 & m2 & $+0.003\,[-0.023,+0.030]$ & 0.366 & 120 & 0.95 & $+0.043$\,$[-0.02,+0.11]$ \\
sweep\_lda-allq50\_clamp\_q50 & m2 & $+0.003\,[-0.020,+0.030]$ & 0.365 & 123 & 0.96 & $+0.023$\,$[-0.03,+0.08]$ \\
gated ablation & m2 & $+0.030\,[-0.013,+0.073]$ & 0.337 & 115 & 0.89 & $+0.119$\,$[+0.03,+0.21]$ \\
sweep\_ocwcr-crq50\_clamp\_q30 & m2 & $+0.020\,[-0.020,+0.063]$ & 0.347 & 117 & 0.91 & $+0.129$\,$[+0.05,+0.21]$ \\
gated clamp $q_{50}$ (sweep) & m2 & $+0.023\,[-0.020,+0.067]$ & 0.355 & 119 & 0.90 & $+0.101$\,$[+0.02,+0.18]$ \\
Linear-AcT (tuned) & m2 & $+0.007\,[-0.027,+0.040]$ & 0.374 & 124 & 0.92 & $+0.052$\,$[-0.02,+0.12]$ \\
additive CAA (tuned) & m2 & $+0.007\,[-0.043,+0.053]$ & 0.330 & 117 & 0.83 & $+0.108$\,$[+0.02,+0.20]$ \\
CAST-style (tuned) & m2 & $+0.007\,[-0.033,+0.047]$ & 0.296 & 103 & 0.87 & $+0.152$\,$[+0.06,+0.25]$ \\
gated full-rank transport (tuned, ours) & m2 & $+0.000\,[-0.030,+0.030]$ & 0.369 & 123 & 0.92 & $+0.036$\,$[-0.03,+0.11]$ \\
MiMiC (tuned) & m2 & $+0.030\,[-0.013,+0.073]$ & 0.352 & 117 & 0.91 & $+0.123$\,$[+0.04,+0.21]$ \\
sequential gate (tuned, ours) & m2 & $+0.007\,[-0.020,+0.033]$ & 0.365 & 122 & 0.95 & $+0.033$\,$[-0.02,+0.09]$ \\
gated local action (tuned, ours) & m2 & $+0.003\,[-0.033,+0.043]$ & 0.357 & 121 & 0.90 & $+0.053$\,$[-0.02,+0.13]$ \\
gated $k$-dim transport (tuned, ours) & m2 & $-0.007\,[-0.033,+0.017]$ & 0.379 & 125 & 0.95 & $+0.011$\,$[-0.04,+0.07]$ \\
tuned\_ourspost & m2 & $-0.003\,[-0.017,+0.007]$ & 0.374 & 125 & 0.99 & $-0.004$\,$[-0.03,+0.02]$ \\
score-proportional dose (tuned, ours) & m2 & $-0.013\,[-0.057,+0.033]$ & 0.371 & 126 & 0.85 & $+0.046$\,$[-0.04,+0.13]$ \\
gated token-local clamp (tuned, ours) & m2 & $-0.003\,[-0.040,+0.030]$ & 0.360 & 122 & 0.90 & $+0.053$\,$[-0.02,+0.13]$ \\
\hline

\end{tabular}}
\caption{MMLU: all conditions, both readouts.}
\end{table}

\begin{table}[!t]\scriptsize\setlength{\tabcolsep}{2.6pt}
\centering
\resizebox{\linewidth}{!}{%
\begin{tabular}{l|c|cc|cc|c}
\hline
condition & r/o & $\Delta$acc [CI] & ECE & \ocw{} & cr\_keep & Sel [CI] \\
\hline
baseline & m4 & --- & 0.092 & 37 & --- & --- \\
additive $-0.75$ & m4 & $-0.123\,[-0.187,-0.060]$ & 0.265 & 0 & 0.00 & $+0.000$\,$[+0.00,+0.00]$ \\
clamp $q_{50}$ & m4 & $+0.037\,[-0.013,+0.087]$ & 0.116 & 41 & 0.82 & $+0.197$\,$[+0.03,+0.36]$ \\
gated clamp $q_{50}$ & m4 & $+0.027\,[-0.010,+0.063]$ & 0.124 & 41 & 0.89 & $+0.137$\,$[-0.01,+0.29]$ \\
quantile-OT non-\ocw{} & m4 & $-0.027\,[-0.060,+0.007]$ & 0.073 & 36 & 0.91 & $+0.095$\,$[-0.03,+0.24]$ \\
NP-gate$\times$BW all\_q50 & m4 & $-0.017\,[-0.040,+0.007]$ & 0.090 & 38 & 0.94 & $+0.052$\,$[-0.05,+0.17]$ \\
subspace-BW calibrated & m4 & $-0.007\,[-0.043,+0.030]$ & 0.127 & 40 & 0.91 & $+0.048$\,$[-0.07,+0.18]$ \\
sweep\_ocwcr-crq50\_ablate & m4 & $-0.003\,[-0.040,+0.033]$ & 0.118 & 38 & 0.84 & $+0.222$\,$[+0.05,+0.39]$ \\
baseline & m2 & --- & 0.229 & 75 & --- & --- \\
additive $-0.75$ & m2 & $-0.200\,[-0.260,-0.140]$ & 0.126 & 90 & 0.56 & $+0.016$\,$[-0.11,+0.15]$ \\
clamp $q_{50}$ & m2 & $-0.020\,[-0.060,+0.020]$ & 0.247 & 81 & 0.90 & $+0.128$\,$[+0.03,+0.23]$ \\
gated clamp $q_{50}$ & m2 & $-0.027\,[-0.063,+0.003]$ & 0.255 & 83 & 0.92 & $+0.057$\,$[-0.02,+0.14]$ \\
quantile-OT non-\ocw{} & m2 & $+0.003\,[-0.017,+0.023]$ & 0.223 & 71 & 0.99 & $+0.080$\,$[+0.02,+0.15]$ \\
NP-gate$\times$BW all\_q50 & m2 & $-0.003\,[-0.020,+0.013]$ & 0.231 & 74 & 0.99 & $+0.040$\,$[-0.01,+0.10]$ \\
subspace-BW calibrated & m2 & $-0.017\,[-0.040,+0.003]$ & 0.250 & 80 & 0.97 & $+0.022$\,$[-0.03,+0.08]$ \\
sweep\_ocwcr-crq50\_ablate & m2 & $-0.007\,[-0.040,+0.023]$ & 0.234 & 77 & 0.94 & $+0.097$\,$[+0.01,+0.19]$ \\
\hline

\end{tabular}}
\caption{ARC-Challenge (transfer): all conditions, both readouts.}
\end{table}

\begin{table}[!t]\scriptsize\setlength{\tabcolsep}{2.6pt}
\centering
\resizebox{\linewidth}{!}{%
\begin{tabular}{l|c|cc|cc|c}
\hline
condition & r/o & $\Delta$acc [CI] & ECE & \ocw{} & cr\_keep & Sel [CI] \\
\hline
baseline & m4 & --- & 0.258 & 6 & --- & --- \\
Linear-AcT $\lambda{=}0.5$ & m4 & $+0.043\,[+0.000,+0.087]$ & 0.297 & 5 & 0.64 & $-0.030$\,$[-0.50,+0.56]$ \\
Linear-AcT $\lambda{=}1.0$ & m4 & $+0.023\,[-0.027,+0.070]$ & 0.269 & 6 & 0.73 & $+0.227$\,$[-0.31,+0.76]$ \\
CAST-style gate$+$add & m4 & $-0.240\,[-0.313,-0.170]$ & 0.055 & 0 & 0.00 & $+0.000$\,$[+0.00,+0.00]$ \\
ablation & m4 & $-0.033\,[-0.090,+0.023]$ & 0.245 & 9 & 0.18 & $-0.151$\,$[-0.67,+0.29]$ \\
additive $-0.75$ & m4 & $-0.267\,[-0.337,-0.200]$ & 0.033 & 0 & 0.00 & $+0.000$\,$[+0.00,+0.00]$ \\
clamp $q_{30}$ & m4 & $-0.003\,[-0.067,+0.057]$ & 0.260 & 8 & 0.36 & $-0.136$\,$[-0.67,+0.43]$ \\
clamp $q_{50}$ & m4 & $+0.027\,[-0.027,+0.077]$ & 0.289 & 8 & 0.45 & $+0.121$\,$[-0.46,+0.62]$ \\
clamp $q_{70}$ & m4 & $+0.033\,[-0.017,+0.087]$ & 0.300 & 10 & 0.55 & $+0.045$\,$[-0.50,+0.60]$ \\
gated clamp $q_{50}$ & m4 & $+0.030\,[-0.020,+0.080]$ & 0.287 & 7 & 0.55 & $+0.212$\,$[-0.36,+0.70]$ \\
gated OT $q_{50}$ & m4 & $+0.040\,[-0.007,+0.087]$ & 0.297 & 5 & 0.55 & $+0.045$\,$[-0.50,+0.62]$ \\
gated clamp $q_{70}$ & m4 & $+0.017\,[-0.030,+0.063]$ & 0.279 & 7 & 0.55 & $+0.212$\,$[-0.33,+0.71]$ \\
gated OT $q_{70}$ & m4 & $+0.010\,[-0.030,+0.047]$ & 0.268 & 5 & 0.64 & $-0.030$\,$[-0.50,+0.55]$ \\
Gaussian-OT cal. & m4 & $+0.020\,[-0.027,+0.063]$ & 0.280 & 5 & 0.73 & $+0.061$\,$[-0.43,+0.62]$ \\
quantile-OT cal. & m4 & $+0.047\,[+0.003,+0.093]$ & 0.295 & 5 & 0.64 & $+0.136$\,$[-0.42,+0.70]$ \\
quantile transport & m4 & $+0.003\,[-0.037,+0.043]$ & 0.257 & 5 & 0.73 & $+0.227$\,$[-0.33,+0.80]$ \\
MiMiC full-space affine & m4 & $+0.063\,[+0.007,+0.120]$ & 0.281 & 11 & 0.45 & $-0.045$\,$[-0.60,+0.50]$ \\
sequential-gate ablation & m4 & $+0.000\,[+0.000,+0.000]$ & 0.258 & 6 & 1.00 & $+0.000$\,$[+0.00,+0.00]$ \\
probegate-all\_q60\_ablate & m4 & $+0.007\,[-0.013,+0.027]$ & 0.262 & 6 & 0.91 & $+0.076$\,$[-0.25,+0.50]$ \\
probegate-all\_q60\_clamp\_q30 & m4 & $+0.000\,[-0.017,+0.017]$ & 0.256 & 8 & 1.00 & $+0.000$\,$[+0.00,+0.00]$ \\
probegate-all\_q60\_otq\_cal & m4 & $+0.010\,[-0.003,+0.023]$ & 0.257 & 6 & 1.00 & $+0.167$\,$[+0.00,+0.50]$ \\
probegate-all\_q75\_ablate & m4 & $+0.007\,[-0.007,+0.020]$ & 0.262 & 6 & 0.91 & $+0.076$\,$[-0.25,+0.50]$ \\
probegate-all\_q75\_clamp\_q30 & m4 & $+0.003\,[-0.007,+0.013]$ & 0.262 & 6 & 1.00 & $+0.167$\,$[+0.00,+0.50]$ \\
probegate-all\_q75\_otq\_cal & m4 & $+0.000\,[-0.010,+0.010]$ & 0.257 & 6 & 1.00 & $+0.167$\,$[+0.00,+0.50]$ \\
prompting & m4 & $-0.030\,[-0.090,+0.033]$ & 0.195 & 15 & 0.27 & $+0.106$\,$[-0.36,+0.50]$ \\
LDA-gated ablation & m4 & $+0.000\,[+0.000,+0.000]$ & 0.258 & 6 & 1.00 & $+0.000$\,$[+0.00,+0.00]$ \\
sweep\_lda-allq50\_clamp\_q30 & m4 & $+0.000\,[+0.000,+0.000]$ & 0.258 & 6 & 1.00 & $+0.000$\,$[+0.00,+0.00]$ \\
sweep\_lda-allq50\_clamp\_q50 & m4 & $+0.000\,[+0.000,+0.000]$ & 0.258 & 6 & 1.00 & $+0.000$\,$[+0.00,+0.00]$ \\
gated ablation & m4 & $-0.050\,[-0.107,+0.010]$ & 0.235 & 11 & 0.18 & $-0.318$\,$[-0.86,+0.22]$ \\
sweep\_ocwcr-crq50\_clamp\_q30 & m4 & $-0.007\,[-0.067,+0.047]$ & 0.267 & 10 & 0.36 & $-0.303$\,$[-0.79,+0.25]$ \\
gated clamp $q_{50}$ (sweep) & m4 & $+0.030\,[-0.020,+0.080]$ & 0.287 & 7 & 0.55 & $+0.212$\,$[-0.36,+0.70]$ \\
Linear-AcT (tuned) & m4 & $+0.023\,[-0.027,+0.070]$ & 0.269 & 6 & 0.73 & $+0.227$\,$[-0.31,+0.76]$ \\
additive CAA (tuned) & m4 & $-0.030\,[-0.093,+0.027]$ & 0.243 & 1 & 0.00 & $+0.000$\,$[+0.00,+0.00]$ \\
CAST-style (tuned) & m4 & $-0.063\,[-0.120,-0.010]$ & 0.218 & 7 & 0.27 & $+0.106$\,$[-0.33,+0.50]$ \\
gated full-rank transport (tuned, ours) & m4 & $+0.000\,[+0.000,+0.000]$ & 0.258 & 6 & 1.00 & $+0.000$\,$[+0.00,+0.00]$ \\
MiMiC (tuned) & m4 & $-0.013\,[-0.070,+0.040]$ & 0.233 & 11 & 0.55 & $+0.045$\,$[-0.47,+0.59]$ \\
sequential gate (tuned, ours) & m4 & $+0.000\,[+0.000,+0.000]$ & 0.258 & 6 & 1.00 & $+0.000$\,$[+0.00,+0.00]$ \\
gated ablation (tuned, ours) & m4 & $+0.000\,[+0.000,+0.000]$ & 0.258 & 6 & 1.00 & $+0.000$\,$[+0.00,+0.00]$ \\
gated $k$-dim transport (tuned, ours) & m4 & $+0.000\,[+0.000,+0.000]$ & 0.258 & 6 & 1.00 & $+0.000$\,$[+0.00,+0.00]$ \\
tuned\_ourspost & m4 & $+0.000\,[+0.000,+0.000]$ & 0.258 & 6 & 1.00 & $+0.000$\,$[+0.00,+0.00]$ \\
score-proportional dose (tuned, ours) & m4 & $+0.003\,[-0.007,+0.013]$ & 0.259 & 6 & 1.00 & $+0.000$\,$[+0.00,+0.00]$ \\
gated token-local clamp (tuned, ours) & m4 & $+0.000\,[+0.000,+0.000]$ & 0.258 & 6 & 1.00 & $+0.000$\,$[+0.00,+0.00]$ \\
baseline & m2 & --- & 0.216 & 87 & --- & --- \\
Linear-AcT $\lambda{=}0.5$ & m2 & $+0.013\,[-0.017,+0.047]$ & 0.204 & 85 & 0.94 & $+0.092$\,$[+0.01,+0.18]$ \\
Linear-AcT $\lambda{=}1.0$ & m2 & $+0.020\,[-0.010,+0.053]$ & 0.202 & 87 & 0.95 & $+0.121$\,$[+0.04,+0.21]$ \\
CAST-style gate$+$add & m2 & $-0.280\,[-0.347,-0.217]$ & 0.226 & 102 & 0.30 & $-0.107$\,$[-0.23,+0.01]$ \\
ablation & m2 & $+0.030\,[-0.020,+0.077]$ & 0.171 & 70 & 0.88 & $+0.335$\,$[+0.22,+0.45]$ \\
additive $-0.75$ & m2 & $-0.290\,[-0.360,-0.220]$ & 0.218 & 105 & 0.26 & $-0.081$\,$[-0.20,+0.03]$ \\
clamp $q_{30}$ & m2 & $+0.040\,[+0.000,+0.073]$ & 0.192 & 71 & 0.94 & $+0.283$\,$[+0.18,+0.38]$ \\
clamp $q_{50}$ & m2 & $+0.013\,[-0.033,+0.057]$ & 0.213 & 74 & 0.89 & $+0.288$\,$[+0.17,+0.40]$ \\
clamp $q_{70}$ & m2 & $+0.030\,[-0.010,+0.070]$ & 0.192 & 69 & 0.94 & $+0.283$\,$[+0.18,+0.39]$ \\
gated clamp $q_{50}$ & m2 & $+0.030\,[-0.010,+0.070]$ & 0.199 & 78 & 0.91 & $+0.222$\,$[+0.12,+0.33]$ \\
gated OT $q_{50}$ & m2 & $+0.023\,[-0.013,+0.060]$ & 0.191 & 76 & 0.94 & $+0.168$\,$[+0.07,+0.26]$ \\
gated clamp $q_{70}$ & m2 & $+0.030\,[-0.010,+0.067]$ & 0.196 & 77 & 0.92 & $+0.181$\,$[+0.09,+0.28]$ \\
gated OT $q_{70}$ & m2 & $+0.017\,[-0.017,+0.050]$ & 0.205 & 82 & 0.93 & $+0.128$\,$[+0.04,+0.22]$ \\
Gaussian-OT cal. & m2 & $-0.010\,[-0.043,+0.027]$ & 0.224 & 91 & 0.92 & $+0.055$\,$[-0.02,+0.14]$ \\
quantile-OT cal. & m2 & $+0.003\,[-0.027,+0.033]$ & 0.201 & 80 & 0.94 & $+0.138$\,$[+0.05,+0.23]$ \\
quantile transport & m2 & $+0.000\,[-0.037,+0.033]$ & 0.212 & 85 & 0.92 & $+0.078$\,$[-0.00,+0.17]$ \\
MiMiC full-space affine & m2 & $+0.030\,[-0.020,+0.080]$ & 0.194 & 77 & 0.87 & $+0.307$\,$[+0.19,+0.42]$ \\
sequential-gate ablation & m2 & $+0.000\,[+0.000,+0.000]$ & 0.216 & 87 & 1.00 & $+0.000$\,$[+0.00,+0.00]$ \\
probegate-all\_q60\_ablate & m2 & $-0.007\,[-0.027,+0.010]$ & 0.220 & 89 & 0.97 & $+0.003$\,$[-0.04,+0.05]$ \\
probegate-all\_q60\_clamp\_q30 & m2 & $+0.003\,[-0.010,+0.017]$ & 0.212 & 86 & 0.98 & $+0.019$\,$[-0.02,+0.07]$ \\
probegate-all\_q60\_otq\_cal & m2 & $+0.003\,[-0.013,+0.020]$ & 0.216 & 86 & 0.98 & $+0.030$\,$[-0.01,+0.08]$ \\
probegate-all\_q75\_ablate & m2 & $-0.003\,[-0.020,+0.010]$ & 0.222 & 91 & 0.98 & $-0.004$\,$[-0.03,+0.03]$ \\
probegate-all\_q75\_clamp\_q30 & m2 & $+0.013\,[+0.003,+0.027]$ & 0.208 & 86 & 0.99 & $+0.006$\,$[-0.01,+0.03]$ \\
probegate-all\_q75\_otq\_cal & m2 & $-0.010\,[-0.023,+0.000]$ & 0.224 & 90 & 0.98 & $-0.004$\,$[-0.03,+0.03]$ \\
prompting & m2 & $+0.010\,[-0.043,+0.063]$ & 0.226 & 88 & 0.83 & $+0.191$\,$[+0.08,+0.31]$ \\
LDA-gated ablation & m2 & $+0.000\,[+0.000,+0.000]$ & 0.216 & 87 & 1.00 & $+0.000$\,$[+0.00,+0.00]$ \\
sweep\_lda-allq50\_clamp\_q30 & m2 & $+0.000\,[+0.000,+0.000]$ & 0.216 & 87 & 1.00 & $+0.000$\,$[+0.00,+0.00]$ \\
sweep\_lda-allq50\_clamp\_q50 & m2 & $+0.000\,[+0.000,+0.000]$ & 0.216 & 87 & 1.00 & $+0.000$\,$[+0.00,+0.00]$ \\
gated ablation & m2 & $+0.013\,[-0.040,+0.063]$ & 0.190 & 76 & 0.84 & $+0.258$\,$[+0.14,+0.37]$ \\
sweep\_ocwcr-crq50\_clamp\_q30 & m2 & $+0.043\,[+0.003,+0.083]$ & 0.179 & 72 & 0.93 & $+0.295$\,$[+0.19,+0.40]$ \\
gated clamp $q_{50}$ (sweep) & m2 & $+0.030\,[-0.010,+0.070]$ & 0.199 & 78 & 0.91 & $+0.222$\,$[+0.12,+0.33]$ \\
Linear-AcT (tuned) & m2 & $+0.020\,[-0.010,+0.053]$ & 0.202 & 87 & 0.95 & $+0.121$\,$[+0.04,+0.21]$ \\
additive CAA (tuned) & m2 & $+0.023\,[-0.030,+0.077]$ & 0.133 & 69 & 0.80 & $+0.343$\,$[+0.23,+0.46]$ \\
CAST-style (tuned) & m2 & $-0.060\,[-0.110,-0.013]$ & 0.204 & 77 & 0.76 & $+0.182$\,$[+0.05,+0.30]$ \\
gated full-rank transport (tuned, ours) & m2 & $+0.000\,[+0.000,+0.000]$ & 0.216 & 87 & 1.00 & $+0.000$\,$[+0.00,+0.00]$ \\
MiMiC (tuned) & m2 & $-0.007\,[-0.053,+0.040]$ & 0.249 & 91 & 0.86 & $+0.147$\,$[+0.04,+0.26]$ \\
sequential gate (tuned, ours) & m2 & $+0.000\,[+0.000,+0.000]$ & 0.216 & 87 & 1.00 & $+0.000$\,$[+0.00,+0.00]$ \\
gated ablation (tuned, ours) & m2 & $+0.000\,[+0.000,+0.000]$ & 0.216 & 87 & 1.00 & $+0.000$\,$[+0.00,+0.00]$ \\
gated $k$-dim transport (tuned, ours) & m2 & $+0.000\,[+0.000,+0.000]$ & 0.216 & 87 & 1.00 & $+0.000$\,$[+0.00,+0.00]$ \\
tuned\_ourspost & m2 & $+0.000\,[+0.000,+0.000]$ & 0.216 & 87 & 1.00 & $+0.000$\,$[+0.00,+0.00]$ \\
score-proportional dose (tuned, ours) & m2 & $+0.000\,[-0.010,+0.010]$ & 0.217 & 85 & 0.99 & $+0.013$\,$[-0.02,+0.05]$ \\
gated token-local clamp (tuned, ours) & m2 & $+0.000\,[+0.000,+0.000]$ & 0.216 & 87 & 1.00 & $+0.000$\,$[+0.00,+0.00]$ \\
\hline

\end{tabular}}
\caption{GSM8K-MCQ (transfer): all conditions, both readouts. Note the tiny
\ocw{}/\crok{} poles (6/12 of 300 at baseline); the model is underconfident
on this domain, selectivity CIs are wide, and transport maps correctly act
near-identity.}
\end{table}

\section{Capability outside the target behavior}
\label{app:capability}

The operators are fitted on multiple-choice reasoning traces, so the question
is what they do to generation that has nothing to do with the target behavior.
Table~\ref{tab:capability} reports code synthesis (HumanEval pass@1, greedy,
$n{=}164$), held-out language modelling (WikiText-103 perplexity, $n{=}300$
passages), and open-ended instruction following (200 AlpacaEval prompts, 384
new tokens): fluency as the perplexity the unsteered model assigns to the
continuation, repetition as distinct-3, and length in words. The last column is
the fraction of prompts on which the sequential gate fires at all.

\begin{table}[!t]\footnotesize\setlength{\tabcolsep}{4pt}
\centering
\resizebox{\linewidth}{!}{%
\begin{tabular}{l|c|c|ccc|c}
\hline
 & HumanEval & WikiText & \multicolumn{3}{c|}{open-ended (AlpacaEval prompts)} & \\
method & pass@1 & ppl & fluency ppl & distinct-3 & win rate & gate fires \\
\hline
unsteered & 0.396 & 23.4 & 3.00 & 0.953 & -- & -- \\
additive $-0.75$ \cite{rimsky2024caa} & 0.128 & 47.3 & 3.36 & 0.480 & 0.043 & -- \\
additive, tuned dose \cite{rimsky2024caa} & 0.408 & 23.6 & 3.14 & 0.922 & 0.495 & -- \\
MiMiC \cite{singh2024mimic} & 0.396 & 24.2 & 3.01 & 0.949 & 0.512 & -- \\
MiMiC, tuned layer \cite{singh2024mimic} & 0.384 & 24.6 & 3.02 & 0.948 & 0.458 & -- \\
Linear-AcT \cite{rodriguez2025act} & 0.390 & 23.4 & 2.98 & 0.950 & 0.505 & -- \\
clamp $q_{50}$ (ours) & 0.390 & 23.1 & 2.97 & 0.946 & 0.507 & -- \\
quantile-OT (ours) & 0.396 & 23.3 & 2.99 & 0.951 & 0.492 & -- \\
online-gated ablation (ours) & 0.396 & n/a & 3.01 & 0.946 & 0.477 & 0.680 \\
gated full-rank transport (ours) & 0.396 & n/a & 3.02 & 0.954 & 0.470 & 0.586 \\
\hline

\end{tabular}}
\caption{Capability preservation away from the fitted behavior. The
unconditional rows keep the action on for every token, so they bound what the
gated operator can cost. WikiText is scored teacher-forced, which gives a
sequential gate no generated prefix to read, so that cell is not applicable.
The win rate is order-consistent against the unsteered answer, judged by
Qwen2.5-7B-Instruct in both presentation orders with ties counted as half, so
$0.5$ means indistinguishable.}
\label{tab:capability}
\end{table}

\section{Do steered activations stay on the manifold?}
\label{app:manifold}

We teacher-force the unsteered traces, apply each operator to the layer-14
residuals, and measure how far the result travels. Table~\ref{tab:manifold}
reports the relative displacement $\lVert h'-h\rVert/\lVert h\rVert$, the 99th
percentile of the norm ratio $\lVert h'\rVert/\lVert h\rVert$, and the distance
to the nearest natural activation in a bank of 4{,}096 unsteered residuals
relative to the unsteered nearest-neighbour distance. We include the transfer
domains ARC and GSM8K because that is where a map fitted on MMLU could be
expected to explode.

One negative result is worth stating rather than hiding. A \emph{full-space}
Mahalanobis distance is useless here: an operator that edits a $k$-dimensional
projection leaves $d-k$ coordinates untouched, which dilute the statistic until
it moves by under $1\%$ for every operator in the table, including the one that
destroys generation. The same dilution has been reported
independently~\cite{vogels2025ids}. We therefore condition instead. Writing
$h=(q,h_\perp)$ with $q=Vh$, the law of $q$ given the 64 leading principal
directions of the complement is Gaussian under the layer's fitted moments, and
its squared Mahalanobis distance is $\chi^2_k$ under the null, so the natural
readout is the fraction of tokens outside the $99\%$ conditional ellipsoid,
against an expected $1\%$. We also report the behavioral check the refusal
literature uses as a selection criterion, the mean KL between the unsteered and
steered next-token distributions with $0.1$ as the accepted
threshold~\cite{arditi2024refusal}, measured against a norm-matched random
direction so that the number is interpretable~\cite{turner2023actadd}. We claim
only that a projection-local edit lands closer to the states the model produces
by itself, not that it lands on them: steered states are provably unreachable
from any discrete prompt~\cite{mishra2026nonsurjective}.

\begin{table}[!t]\footnotesize\setlength{\tabcolsep}{3pt}
\centering
\resizebox{\linewidth}{!}{%
\begin{tabular}{l|cccc|cccc|cccc}
\hline
 & \multicolumn{4}{c|}{MMLU} & \multicolumn{4}{c|}{ARC} & \multicolumn{4}{c}{GSM8K} \\
operator & disp. & NN & out\% & KL & disp. & NN & out\% & KL & disp. & NN & out\% & KL \\
\hline
unsteered & --- & 1.00 & 0.9 & 0.000 & --- & 1.00 & 1.2 & 0.000 & --- & 1.00 & 2.9 & 0.000 \\
random direction, norm-matched & 0.663 & 1.90 & 1.6 & 0.445 & 0.644 & 2.27 & 1.7 & 0.233 & 0.749 & 2.06 & 0.9 & 0.149 \\
additive $-0.75$ \cite{rimsky2024caa} & 0.663 & 1.88 & 100.0 & 0.032 & 0.644 & 2.25 & 100.0 & 0.008 & 0.749 & 2.04 & 100.0 & 0.012 \\
additive, tuned dose \cite{rimsky2024caa} & 0.221 & 1.22 & 69.9 & 0.005 & 0.215 & 1.33 & 40.3 & 0.001 & -- & -- & -- & -- \\
MiMiC \cite{singh2024mimic} & 0.082 & 1.06 & 1.3 & 0.001 & 0.074 & 1.08 & 1.8 & 0.000 & 0.097 & 1.06 & 4.5 & 0.001 \\
Linear-AcT \cite{rodriguez2025act} & 0.012 & 1.01 & 1.0 & 0.000 & 0.011 & 1.01 & 1.5 & 0.000 & 0.012 & 1.01 & 3.5 & 0.000 \\
clamp $q_{50}$ (ours) & 0.034 & 1.03 & 1.3 & 0.006 & 0.046 & 1.06 & 1.1 & 0.001 & 0.003 & 1.00 & 2.9 & 0.001 \\
quantile transport (ours) & 0.004 & 1.00 & 1.0 & 0.000 & 0.003 & 1.00 & 1.3 & 0.000 & 0.007 & 1.00 & 4.0 & 0.000 \\
ablation \cite{arditi2024refusal} & 0.144 & 1.15 & 29.5 & 0.010 & 0.178 & 1.28 & 28.4 & 0.001 & 0.062 & 1.05 & 10.1 & 0.001 \\
\hline

\end{tabular}}
\caption{On-manifold diagnostics at layer 14. NN is the mean
nearest-neighbour distance to a bank of natural activations relative to the
unsteered value, so $1.00$ means the edit lands as close to the model's own
states as the states it replaced; out\% is the fraction of tokens outside the
$99\%$ conditional ellipsoid, against $1\%$ expected; KL is the mean divergence
of the next-token distribution, with $0.1$ the accepted
threshold~\cite{arditi2024refusal}.}
\label{tab:manifold}
\end{table}

\section{End-to-end serving cost}
\label{app:serving}

Table~\ref{tab:serving} measures the whole 300-question workload on one H100,
so the post-hoc gate is charged for its scoring pass and for regenerating what
it flags.

\begin{table}[!t]\footnotesize\setlength{\tabcolsep}{5pt}
\centering
\begin{tabular}{l|c|cc}
\hline
method & passes & seconds & wall clock \\
\hline
unsteered & 1 & 43.5 & 1.00$\times$ \\
additive $-0.75$ \cite{rimsky2024caa} & 1 & 45.2 & 1.04$\times$ \\
ablation \cite{arditi2024refusal} & 1 & 44.6 & 1.02$\times$ \\
clamp $q_{50}$ (ours) & 1 & 41.4 & 0.95$\times$ \\
quantile-OT (ours) & 1 & 44.8 & 1.03$\times$ \\
MiMiC \cite{singh2024mimic} & 1 & 42.9 & 0.98$\times$ \\
Linear-AcT \cite{rodriguez2025act} & 1 & 40.9 & 0.94$\times$ \\
gated ablation, post-hoc gate (ours) (223/300 regenerated) & 2 & 85.8 & 1.97$\times$ \\
gated ablation, sequential gate (ours) & 1 & 45.5 & 1.05$\times$ \\
\hline

\end{tabular}
\caption{End-to-end wall clock, MMLU $n{=}300$, batch 64, one H100.}
\label{tab:serving}
\end{table}

\section{Definitions used in the text}
\label{app:defs}

\textbf{Hedge markers.} The lexicon counted in Sec.~\ref{sec:results} and
Appendix~\ref{app:samples} is fixed in advance and case-insensitive:
\emph{maybe, perhaps, might, possibly, however, wait, actually, hmm, unsure,
unclear, guess, seems, probably, could, uncertain, confused}. The matching
certainty lexicon is \emph{clearly, obviously, definitely, certainly, surely,
undoubtedly, evidently, must, certain, sure}. A trace's hedge count is the
number of token matches; rates are per trace.

\textbf{Calibrated zone.} The green ellipse in Fig.~\ref{fig:plane} is the
$2\sigma$ contour of the empirical Gaussian fitted to the trace-mean
projections of the \emph{calibrated} extraction-split states
(\crok{} and nonconfident-wrong), i.e.\ the states the transport target is
built from. It is a visual reference, not a decision boundary.

\textbf{Norm-matched spherical action.} The comparison action in
Table~\ref{tab:postreview} is a spherical steering operator: activations are
rotated along the fixed-azimuth shortest geodesic towards a unit prototype at
constant norm, with the decision taken by an antipodal von~Mises--Fisher
confidence at concentration $\kappa{=}20$. Its nine-cell grid
($\alpha\in\{0.6,0.9,1\}$, $\beta\in\{0.6,0.9,0.99\}$) is selected on the
observer split, and the reported comparison uses the selected cell, so the two
conditions differ in the direction and the decision rule and not in how far
they move activations. The five \emph{isotropic} controls replace the detector
direction with a random unit direction at exactly the measured update norm;
the \emph{unprompted} control refits the direction at last-token positions
without the evaluation instruction. All are input-independent given the fit.

\textbf{Condition names.} The exhaustive grids in
Appendix~\ref{app:fulltables} carry the run identifiers; the main text uses
short names. \texttt{m4\_conf\_*} is an action on the confidence axis;
\texttt{sweep\_ocwcr-crq50\_ablate} is gated ablation (gate fitted on the
\ocw{}-versus-\crok{} contrast, threshold at the \crok{} median);
\texttt{sweep\_lda-allq50\_ablate} is LDA-gated ablation (the two-dimensional
LDA head of Sec.~\ref{sec:method}, threshold at the median of all extraction
scores); \texttt{online\_d0p05\_q50\_L16\_ablate} is its sequential
single-pass form at $\delta{=}0.05$; \texttt{probegate-*} gates on a full-$d$
logistic probe; \texttt{cast\_gate-cr\_q50\_add-0.75} is the CAST-style
reimplementation; \texttt{act\_l05}/\texttt{act\_l10} are Linear-AcT at
$\lambda{=}0.5$/$1.0$; \texttt{mimic\_full} is MiMiC;
\texttt{subspace-BW} is ungated subspace Bures--Wasserstein transport and
\texttt{NP-gate}$\times$\texttt{BW} its gated form (\texttt{subbw} and
\texttt{npbw} in the run identifiers); \texttt{self\_qXX} marks a threshold
re-anchored, label-free, to the target domain's own score quantile; and
\texttt{-think} marks steering confined to the \texttt{<think>} phase.

\selectlanguage{russian}
% Times (ptm) has no T2A variant in this distribution — render the Russian
% back-matter in Computer Modern (LH fonts).
\FloatBarrier
\begingroup\renewcommand{\rmdefault}{cmr}\normalfont
\begin{center}
\textbf{М.~Каширский, Д.~Сергеев, Т.~Тер-Ованнисян, И.~Пионтковская. Проекционно-транспортный стиринг: условное управление поведением рассуждений языковых моделей}
\end{center}

\begin{abstract}
Стиринг активаций обещает управление поведением без дообучения, однако
глобальный сдвиг остаточного потока подавляет полезное и вредное поведение
разом. Мы рассматриваем стиринг как задачу оптимального транспорта в
низкоразмерном поведенческом подпространстве и отделяем решение о
вмешательстве от самого вмешательства. Доказано, что транспортная карта между
распределениями проекций --- единственное вмешательство минимального смещения,
локальное в этом подпространстве; что селективность любой функции переноса на
оси ограничена разделимостью подпопуляций вдоль неё; и что оптимальный по
Нейману--Пирсону гейт в композиции с картой оптимален совместно в своём
классе. Последовательная форма гейта принимает решение по префиксу, поэтому
оператор работает за один проход. На сверхуверенности рассуждений
\texttt{gemma-2-2b-it} стандартная аддитивная правка устраняет все
уверенно-неверные ответы только ценой уверенно-верных, а измеренное
покомпонентное значение $d=0{,}049$ между этими подпопуляциями делает такой
исход геометрической неизбежностью, а не следствием неудачной настройки. На
пяти переотобранных подмножествах MMLU условный оператор достигает
селективности $+0{,}279$ без потери точности, переносится на
ARC-Challenge и GSM8K без перенастройки и воспроизводится при увеличении
масштаба в $4{,}5$ раза. При равном бюджете настройки у всех методов он
сохраняет эту селективность при смене вопросов, тогда как настроенная
полноранговая карта теряет две трети своей. Он также сохраняет способности
там, где аддитивная правка их разрушает, оставаясь на уровне невмешательства
на HumanEval, WikiText-103 и открытых генерациях, а отложенная проверка на
OLMo-2/TruthfulQA подтверждает конструкцию детекторной оси за пределами
Gemma.
\end{abstract}

\keywords{стиринг активаций \and оптимальный транспорт \and калибровка \and сверхуверенность \and интерпретируемость.}
\endgroup

\end{document}
